%==============================================================================
%  Joint Transport--Relay Co-Design for UAV Emergency Logistics in Mountainous
%  Terrain:  Handover-Aware Feasibility, Channel-Model Sensitivity, and the
%  Price of Partitioning
%
%  English SCI manuscript (draft).  All numbers are macros generated by
%  code/make_numbers_en.py from the result files; see numbers_en_sources.txt
%  for the provenance of every macro.
%==============================================================================
\documentclass[11pt]{article}

\usepackage[margin=2.4cm]{geometry}
\usepackage{amsmath,amssymb,amsthm}
\usepackage{booktabs}
\usepackage{graphicx}
\usepackage{siunitx}
\usepackage{enumitem}
\usepackage[hidelinks]{hyperref}
\usepackage{xcolor}
\usepackage{caption}
\usepackage{algorithm}
\usepackage{algpseudocode}
\algrenewcommand\algorithmicrequire{\textbf{Input:}}
\algrenewcommand\algorithmicensure{\textbf{Output:}}
\floatname{algorithm}{Algorithm}
\captionsetup{font=small}

% ---- \texttt 断行：正文引用 code/make_numbers_en.py 这类长路径，
% IEEEtran 双栏窄栏下无法断行会 overfull。改为每个字符后允许换行。
% 必须先用 \makeatletter 定义辅助宏、再 \renewcommand，否则展开时未定义。
\makeatletter
\def\dsh@tt#1{\dsh@ttloop#1\dsh@end}
\def\dsh@ttloop#1{\ifx#1\dsh@end\else\dsh@ttemit{#1}\expandafter\dsh@ttloop\fi}
\def\dsh@ttemit#1{#1\penalty\z@\hskip\z@ plus 0.02em\relax}
\renewcommand{\texttt}[1]{{\ttfamily\hyphenchar\font=-1 \dsh@tt{#1}}}
\makeatother

% ---- 表格自适应列宽 ---------------------------------------------------------
% 同一个 body 要能在三种版式下编译：单栏 article / elsarticle preprint，
% 以及 IEEEtran 的双栏窄栏。表格按单栏宽度排就会在双栏里溢出（右列被裁掉），
% 所以所有表格统一走 \fitwide：把表格缩放到当前版心的宽度。
%   * 单栏版式：表格本身通常窄于版心，缩放系数 > 1 会被 \resizebox 忽略（只缩不放时
%     用 \fitwide 的 \ifdim 分支），因此不会把小表拉大；
%   * 双栏版式：在 table* 里 \linewidth = \textwidth（跨双栏），表格得到足够宽度。
\newsavebox{\wearfitbox}
\newcommand{\fitwide}[1]{%
  \sbox{\wearfitbox}{#1}%
  \ifdim\wd\wearfitbox>\linewidth
    \resizebox{\linewidth}{!}{\usebox{\wearfitbox}}%
  \else
    \usebox{\wearfitbox}%
  \fi
}



% ---- 结果文件缺失时给出显式提示，而不是静默留空 ----------------------------
% 依序尝试：调用时给出的路径 → 当前目录 → paper_en/ → 上一级的 paper_en/ → 上一级。
% 这样在三种布局下都能编：工作区（在 D题/ 下编 paper_en/main.tex）、
% paper_en/ 内直接编译、以及交付包（06_英文SCI稿/ 内编译）。
\newcommand{\InputRes}[1]{%
  \IfFileExists{#1}{\input{#1}}{%
    \IfFileExists{paper_en/#1}{\input{paper_en/#1}}{%
      \IfFileExists{../paper_en/#1}{\input{../paper_en/#1}}{%
        \IfFileExists{../#1}{\input{../#1}}{%
          \typeout{== WARNING: resource #1 NOT FOUND ==}%
          \par\noindent\fbox{\parbox{0.9\linewidth}{\centering Missing resource
          \texttt{#1}; run \texttt{python code/make\_numbers\_en.py}}%
          }\par
        }%
      }%
    }%
  }%
}
% ---- 插图：优先用矢量 PDF，回退到 300 dpi PNG -------------------------------
% \Dfig[宽度]{不带扩展名的图名}{图注}{label}
% 为什么带扩展名解析：数据图由 make_figs_en.py / make_fig_theory.py 以
% **矢量 PDF + 300 dpi PNG** 双份导出。双栏排版把图缩到约 8.8 cm 宽，
% 只有矢量版本在任意缩放下锐利；而 DEM 晕渲底图本身是栅格，转矢量只会
% 撑大 PDF 不增清晰度，因此**只有 PNG**，这里靠回退逻辑自动取用。
\newcommand{\Dfigfile}[3]{%
  \IfFileExists{#3#2.pdf}{\includegraphics[width=#1\linewidth]{#3#2.pdf}}{\IfFileExists{#3#2.png}{\includegraphics[width=#1\linewidth]{#3#2.png}}{\IfFileExists{#3#2.eps}{\includegraphics[width=#1\linewidth]{#3#2.eps}}{\fbox{\parbox{0.7\linewidth}{\centering Figure \texttt{#2} not yet generated; run \texttt{python code/make\_figs\_en.py}}}}}}%
}
% 插图入口：\Dfig[宽度]{图名（不带扩展名）}{图注}{label}
% 图名一律省略扩展名——由 \Dfigfile 按 pdf → png → eps 的顺序挑，
% 这样"数据图用矢量、DEM 晕渲用 300 dpi 栅格"这件事由文件是否在场决定，
% 正文不需要知道哪张是矢量。
\newcommand{\Dfig}[4][0.92]{%
  \begin{figure}[htbp]\centering
    \Dfigfile{#1}{#2}{figs/}\caption{#3}\label{#4}%
  \end{figure}%
}
% 跨栏版：这些图都是 1×2 / 1×3 的多子图，在 IEEEtran 里必须跨双栏
% （figure*）才放得下；否则会被压到单栏 3.5 in 内，图内 8 pt 的字缩成 ~4 pt。
% 图本身按 7.2 in 画布、8 pt 基准字号生成，与跨栏显示宽度 1:1，不再二次缩小
% —— 这才是"图内文字发虚"的真正修法（矢量图不存在分辨率问题，只有缩放比问题）。
\newcommand{\Dfigwide}[4][0.98]{%
  \begin{figure*}[htbp]\centering
    \Dfigfile{#1}{#2}{figs/}\caption{#3}\label{#4}%
  \end{figure*}%
}

% ---- 标题统一：同一字体族、不用全大写 ----------------------------------------
% 三种版式对 \section 的默认处理不一样：
%   * article（中立版）：\Large\bfseries，本来就保留大小写；
%   * IEEEtran：\normalfont\normalsize\centering\scshape —— 小体大写，
%     于是 "System model" 印成 "SYSTEM MODEL"；
%   * elsarticle：\large\bfseries\scshape 或默认 article 样式。
% 结果是"章标题全大写、节标题正常大小写"，两套样式并存；审稿时看着像格式不统一。
% 这里统一为：同一字体族（正文罗马体）+ 同一字重（粗体）+ 保留**句子大小写**，
% 三级标题（section / subsection / subsubsection）保持一致，不再出现全大写。
% 位置参数沿用各类的默认值，避免破坏正文与标题的间距。
\makeatletter
\providecommand{\dsh@section@spacing}{%
  {1\baselineskip plus 0.25\baselineskip minus 0.25\baselineskip}%
  {0.7ex plus 1ex minus 0ex}}
\@ifclassloaded{IEEEtran}{%
  \def\section{\@startsection{section}{1}{\z@}%
    {1.0\baselineskip plus 0.25\baselineskip minus 0.25\baselineskip}%
    {0.4\baselineskip}%
    {\normalfont\normalsize\centering\bfseries}}%
  \def\subsection{\@startsection{subsection}{2}{\z@}%
    {1.0\baselineskip plus 0.25\baselineskip minus 0.25\baselineskip}%
    {0.3\baselineskip}%
    {\normalfont\normalsize\itshape}}%
  \def\subsubsection{\@startsection{subsubsection}{3}{\z@}%
    {0.8\baselineskip plus 0.2\baselineskip minus 0.2\baselineskip}%
    {0.2\baselineskip}%
    {\normalfont\normalsize\itshape}}%
}{%
  % article（中立版）与 elsarticle：只需去掉 small-caps，字族字重沿用默认
  \@ifundefined{elsarticle@type}{%
    \renewcommand{\section}{\@startsection{section}{1}{\z@}%
      {-3.5ex plus -1ex minus -0.2ex}{2.3ex plus 0.2ex}%
      {\normalfont\Large\bfseries}}%
  }{%
    \renewcommand{\section}{\@startsection{section}{1}{\z@}%
      {-3.5ex plus -1ex minus -0.2ex}{2.3ex plus 0.2ex}%
      {\normalfont\large\bfseries}}%
  }%
}
\makeatother

% ---- 公式自适应字号 ---------------------------------------------------------
% 有几个长行间公式在单栏版式下正好放得下，但在 IEEEtran 的单栏里会溢出
% （\hbox too wide，压到栏外）。\eqfit 用 graphicx 的 \width 直接判断自然宽度：
% 放不下才按比例缩到版心，放得下就维持 100% 原样，因此单栏版式排版不变。
% 注意整条公式只有一个 \resizebox，不要在 \ifdim 分支里各写一份——那样宏内的
% 换行会让 \setbox 少吃掉一个空格并把行号推后一行。
%
% 另外两处细节是排版事故的来源，必须保留：
%   * 开头的 \par 把前一段**收尾**。否则当 \eqfit 紧跟在 \begin{proof} 之后时，
%     "Proof." 之后那一行没有可断处，TeX 会把极少几个词撑满整栏——看起来就是
%     "Proof.        The        cruise        altitude        is" 那种诡异间距。
%   * 结尾的 \par 把 center 段落封口，避免后文与公式之间产生额外空行。
\newcommand{\eqfit}[1]{%
  \par\nobreak
  \begin{center}%
  \resizebox{\ifdim\width>\linewidth\linewidth\else\width\fi}{!}{$\displaystyle #1$}%
  \end{center}%
  \par\nobreak
}

% ---- 浮动体间距：避免图与正文小标题贴得太近（双栏下尤其明显）----------------
\setlength{\textfloatsep}{12pt plus 3pt minus 2pt}
\setlength{\floatsep}{12pt plus 3pt minus 2pt}
\setlength{\intextsep}{12pt plus 3pt minus 2pt}
\renewcommand{\topfraction}{0.92}
\renewcommand{\bottomfraction}{0.85}
\renewcommand{\textfraction}{0.06}
\renewcommand{\floatpagefraction}{0.82}

% ---- 定理类环境：统一为标准版式，标题独立成行 --------------------------------
% amsthm 默认把**标题与正文放在同一段**，于是标题也参与该行的两端对齐。当正文
% 以一条长的不可断行公式开头时（例如 Corollary 5 的
% $D=\sum_k\max(0,w-\mathrm{span}(k-w))$），这一行没有可断处，TeX 只能把标题里的
% 几个空格撑开去填满整栏——双栏版式下就变成
% "Corollary        5        (Deficiency        functional)."
% 这类忽宽忽窄的样子（单栏时栏宽大、不易察觉，双栏时非常明显）。
%
% 修法：把标题与正文之间的分隔（\newtheoremstyle 的第 7 个参数）改成
% "\newline"，让**标题独占一行**（"By name" 风格）。标题行只含标题本身，
% 没有可拉伸的余地；正文另起一行，两端对齐不再受公式长度牵制。
% 标题之后不再加句点，这也与常见的 "By name" 定理样式一致。
\newtheoremstyle{dshplain}%
  {\topsep}{\topsep}%
  {\itshape}%   正文
  {}%           缩进
  {\bfseries}%  标题
  {}%           标题后标点：标题独立成行，不加句点
  {\newline}%   标题与正文之间：换行
  {}%           标题格式
\newtheoremstyle{dshdefn}%
  {\topsep}{\topsep}%
  {\normalfont}% 定义/假设用直立体，与定理类区分
  {}%
  {\bfseries}%
  {}%
  {\newline}%
  {}
\theoremstyle{dshplain}
\newtheorem{theorem}{Theorem}
\newtheorem{proposition}[theorem]{Proposition}
\newtheorem{corollary}[theorem]{Corollary}
\theoremstyle{dshdefn}
\newtheorem{assumption}[theorem]{Assumption}
\newtheorem{definition}[theorem]{Definition}

\newcommand{\tbd}{\texttt{TBD}}
\newcommand{\y}[1]{\textsc{#1}}

% ---- 紧急拉伸：12pt 单栏版式下有成串的行内数学（FSPL 不等式、
% Fresnel-Kirchhoff 参数等），这些数学整体不可断字，偶尔会把一行撑出版心。
% \emergencystretch 只在原本会 overfull 的行上起作用，不影响正常排版。
% 放在这里是为了同时进入三种版式（期刊变体只取"共享段"）。
\setlength{\emergencystretch}{3em}

\InputRes{numbers_en.tex}

\title{\bfseries Joint Transport--Relay Co-Design for UAV Emergency Logistics\\
in Mountainous Terrain:\\ Handover-Aware Feasibility, Channel-Model Sensitivity,\\
and the Price of Partitioning}
\author{Changxuan Cao\textsuperscript{1}, Qiong Li\textsuperscript{1},
Zhifeng Liu\textsuperscript{1,*}, Yulong Peng\textsuperscript{1}, Zexu Ouyang\textsuperscript{1}\\
\small \textsuperscript{1}East China University of Technology (ECUT), Nanchang, Jiangxi, China\\
\small \textsuperscript{*}Corresponding author: \texttt{2023211863@ecut.edu.cn}}
\date{\today}

\begin{document}
\maketitle

%==============================================================================
\begin{abstract}
\noindent
Mountainous disaster relief requires unmanned aerial vehicles (UAVs) to deliver
cargo while staying connected to a fixed gateway that terrain frequently
occludes. We study the resulting \emph{transport--relay co-design} problem and
make five contributions.
(i) We show that the widely used ``relay coverage'' abstraction is
\emph{not executable on our instance}: once the physical cost of \emph{handing
over} between hovering positions is accounted for, the fleet size for which we can
construct a verified schedule grows from $\ENrbNaive$ to $\ENrbGreedy$ vehicles,
and we give a checkable necessary condition for a handover
(Theorem~\ref{thm:handover}) that turns the handover cost into an inequality on a
single scalar \emph{span} function. We report the rate at which that condition is
satisfiable at the median blackout window ($\ENqStarBase\%$ of the demand cells
under M1, $\ENqStarAir\%$ under M2) and show that it does
\emph{not} close the question; we consequently report every fleet-size statement
with its search scope, distinguishing a proven lower bound, a constructed
schedule, and a search that returned no schedule. On the declared candidate set
$C_0$ (refined per cell to the top-$12$ points by remaining coverage) the state
recursion admits no feasible state at $t=\ENgrFailNTwoT$\,s under M1 and
$t=\ENgrFailNTwoTairPP$\,s under M2, and an independent phase-level integer
program shows that the obstruction lives at \emph{instant} rather than phase
granularity; a sampling-free energy bound gives $N^{*}\ge\ENboundEnergy$. That
negative verdict is a statement about the declared finite set, and we say so
wherever it is used: closing the state pruning at that breadth did not terminate
within our compute budget, so we do not claim an unpruned decision.
(ii) We replace the constant occlusion penalty supplied with the benchmark by
the ITU-R P.~526 single-knife-edge diffraction model evaluated on the 30\,m
digital elevation model; this raises the outage demand by $\ENchReqDelta\%$ and
reduces the number of usable relay hover points by $\ENchCandDeltaAbs\%$, and the
coverage-only relaxation needs $\ENrbNaive$ relay while the smallest fleet
for which we can construct a verified schedule rises from $\ENchObsNStar$ to
$\ENchAltNStar$, so the commonly reported ``add one relay''
conclusion must be qualified by the channel convention.
(iii) We close the loop between transport scheduling and relay feasibility with
a threshold-accepting search over departure instants alone, holding the routing,
batching and fleet decisions at their baseline values; the residual handover
deficit was unchanged by that search, which is evidence against a purely
scheduling-driven explanation but not a proof of its absence. The classical
Clarke--Wright savings heuristic, run under the same scheduler, fails from the
opposite direction: it wins the routing layer
($\ENcwSorties$ sorties, $\ENcwEnergy$\,kWh) and violates hard deadlines by
$\ENcwHard$\,s.
(iv) We introduce a controlled terrain family (relief scaling, corridor
elevation and terrain translation) and show that the smallest fleet with a
constructed schedule is a non-monotone function of terrain severity, with no
monotone relation to the occlusion ratio.
(v) Every table and figure value in this paper is generated from the shipped
result files by a single scripted pipeline, with the exceptions listed explicitly
in Section~\ref{sec:reproducible}; every bibliography entry was checked against
its publisher record by DOI.
\end{abstract}

\noindent\textbf{Keywords:} UAV logistics; emergency logistics; drone delivery;
relay handover; fleet sizing; scheduling; disaster relief.

%==============================================================================
\section{Introduction}
\label{sec:intro}

\subsection{Motivation}
Rapid-onset disasters in mountainous regions isolate settlements: roads fail
while the very terrain that isolates them also blocks the line-of-sight (LoS)
links on which UAV command and control depends \cite{sheu2007,caunhye2012,
erdelj,otto2018}. Relief fleets therefore need \emph{both} cargo delivery
and communication relaying, and the two are coupled: a relay must hover where it
can see the transport aircraft, and the transport aircraft can only fly when a
relay sees it. This coupling is what the present paper isolates and quantifies.
The operator's decision variables are the familiar transportation ones --- how
many vehicles, on which routes, dispatched when --- but the feasibility of a
dispatch plan now depends on a communication constraint that switches on and
off over time, so a plan that is optimal for delivery alone may be
\emph{infeasible} to fly. Recent work in this journal on synchronised
truck--drone routing \cite{das2021,bian2024} addresses the routing half of this
coupling; we take the routing layer as given and characterise the half that
decides whether the plan can be executed at all.

\subsection{Gaps in existing treatments}
Three modelling choices in the applied literature are convenient but
consequential, and each is examined here.
\begin{enumerate}[leftmargin=1.5em,itemsep=2pt]
\item \textbf{Instantaneous relocation.} Relays are commonly treated as
  coverage devices that can be repositioned at no cost
  \cite{zeng,mozaffari,zengtraj,alhourani2014,hourani2016}. Repositioning a
  multirotor between two hovering points is not instantaneous: it must climb,
  cruise, descend and re-establish the link, during which it provides
  \emph{no} coverage. Section~\ref{sec:handover} shows this is a first-order
  effect, not a refinement.
\item \textbf{A constant occlusion penalty.} The benchmark instance used here
  supplies a single ``terrain occlusion loss'' in dB, applied whenever the LoS
  is intersected; several applied UAV studies adopt the same convention
  \cite{matolak,alhourani2014}. This treats a grazing ridge and a full mountain alike, and
  -- worse -- declares short blocked links usable. Section~\ref{sec:channel}
  quantifies the resulting bias against the profile-based ITU-R P.~526 model
  \cite{iturp526,deygout,matolak}.
\item \textbf{Frozen transport schedules.} Feasibility of the relay layer is
  usually assessed on a transport schedule optimised for delivery alone, which
  invites the objection that infeasibility is an artefact of decoupling.
  Section~\ref{sec:feedback} answers it by optimising the schedule for relay
  feasibility.
\end{enumerate}

\subsection{Contributions}
\begin{enumerate}[leftmargin=1.8em,itemsep=2pt]
\item \textbf{An executable handover model and an explicit statement of what is
  proven.} We replace the instantaneous-relocation abstraction by the blackout
  window $W(a,b)$ (Definition~\ref{def:blackout}) and prove that a relay schedule
  is feasible only if a scalar \emph{span} function satisfies a checkable
  inequality in every demand cell (Theorem~\ref{thm:handover}). That inequality
  is a necessary condition, not a certificate: we measure how often it is
  satisfiable and find $\ENqStarBase\%$ (M1) and $\ENqStarAir\%$ (M2) of candidate
  pairs at the median blackout window, so it leaves the two-relay question open.
  We therefore classify every fleet-size statement by its evidence --- a proven
  lower bound, a constructed and verified schedule, or a search that returned no
  schedule --- and we state $N^{*}$ as the interval
  $\max(N_{\mathrm{naive}},LB_{\mathrm{handover}})\le N^{*}\le\ENrbGreedy$ rather
  than as a single value (Section~\ref{sec:premium},
  Sections~\ref{sec:premium}--\ref{sec:ilp}). On the declared candidate set
  $C_0$ with a per-cell top-$12$ refinement, the state recursion admits no
  feasible state; the phase-level integer program
  is satisfiable at $N=2$, and the sampling-free energy bound gives
  $N^{*}\ge\ENboundEnergy$.
\item \textbf{A profile-based channel model and its operational consequence.}
  We evaluate ITU-R P.~526 single-knife-edge diffraction on the $30$\,m terrain
  profile and show that it raises outage demand by $\ENchReqDelta\%$ and removes
  $\ENchCandDeltaAbs\%$ of the usable hover points relative to the supplied constant
  penalty, raising the smallest fleet for which we can construct a verified
  schedule from $\ENchObsNStar$ to $\ENchAltNStar$ --- with the outcome unchanged
  across candidate budgets (Section~\ref{sec:channel}).
\item \textbf{A null result on decoupling, and a textbook heuristic that fails
  from the opposite direction.} A feedback loop that re-optimises only the
  departure instants, holding routing, batching and fleet decisions fixed,
  returned no reduction of the handover deficiency (Section~\ref{sec:feedback});
  and when the classical Clarke--Wright savings heuristic \cite{clarke1964} is
  run under the \emph{same} scheduler it wins the routing layer but violates hard
  deadlines by $\ENcwHard$\,s (Section~\ref{sec:ablation}). Both results argue for
  keeping scheduling inside the search loop.
\item \textbf{What the fleet size responds to.} Across a controlled terrain
  family, in which only the elevation field is transformed, the smallest fleet
  with a constructed schedule is \emph{non-monotone} in terrain severity: it is
  uncorrelated with the occlusion ratio (Pearson $\rho=-0.23$, $95\%$ CI
  $[-0.77,+0.52]$) and, among the predictors we computed, most strongly
  associated with temporal fragmentation of the outage demand
  ($\rho=+0.47$, $95\%$ CI $[-0.28,+0.86]$). With nine scenarios --- five of
  which admitted no schedule for $N\le4$ and are treated as censored --- none of
  these associations is statistically established, and we report them as
  descriptive (Section~\ref{sec:terrain}).
\item \textbf{A reproducible artefact.} {\sloppy Every table and figure value in
  this paper is generated from shipped result files by the scripts named below,
  with the exceptions listed in
  Section~\ref{sec:reproducible}; every bibliography entry was checked against its
  publisher record by DOI. The pipeline is one command per stage.}\par
\end{enumerate}

%==============================================================================
\section{Related work}
\label{sec:related}
\paragraph{UAV logistics in disrupted terrain.}
Drone-assisted relief delivery is usually formulated as a vehicle-routing
problem with time windows \cite{solomon1987}, augmented by payload, endurance
and battery-swap constraints; heterogeneous fleets, synchronised truck--drone
operations and multi-trip schedules are standard
\cite{murray,poikonen2017,otto2018,macrina2020,das2021,bian2024,coutinho,boysen,
erdelj,scott2018}. Post-disaster reconnaissance is the complementary question of
\emph{where to look} \cite{oruc2018}, and environmental accounting has been used
to justify drone delivery in the last mile \cite{goodchild2018}. Reviews of civil
UAV optimisation \cite{otto2018,macrina2020} and of the underlying VRP machinery
\cite{tothvigo2002,vidal2020,clarke1964} agree on the standard architecture:
build routes first, then schedule them. Our transport layer follows this line,
but it is deliberately treated as a \emph{subproblem}: the object of study is the
coupling between the delivery schedule and the feasibility of the communication
layer, not the routing objective itself. Section~\ref{sec:ablation} shows, as a
side result, that this standard architecture is exactly where a routing-first
heuristic loses: it wins on route length and sortie count and loses on hard
deadlines.

\paragraph{Relaying and connectivity.}
A parallel literature places ground or aerial relays to restore connectivity,
typically maximising coverage, minimising the number of relays, or maintaining a
connected chain \cite{zeng,mozaffari,zengtraj,zeng2016relay,relay5g,gomez2016}. The
dominant abstraction is a \emph{coverage} relation between a relay position and a
set of terminals, with relocation either free or penalised by a motion cost; for
quasi-static platforms the deployment side of this problem is well understood
\cite{alhourani2014,hourani2016}. Section~\ref{sec:premium} shows that for
\emph{hovering} relays this abstraction is not executable: because a relay
provides no service while transiting between hover points, the fleet size is
governed by a \emph{span} inequality rather than by a coverage count, and on our
instance the coverage abstraction understates the fleet by a factor of
$\ENrbGreedy/\ENrbNaive$. The distinction is not that relocation ``has a cost'' in
the usual motion-energy sense --- it is that the relay is \emph{unavailable} while
relocating, which turns a static covering problem into a continuous-time
scheduling problem.

\paragraph{Air-to-ground propagation in mountains.}
Studies of low-altitude platforms use either free-space loss with an empirical
excess attenuation \cite{matolak}, or a probabilistic LoS model calibrated for
altitude selection \cite{alhourani2014}, or terrain-aware diffraction models;
ITU-R P.~526 \cite{iturp526} and the Deygout construction \cite{deygout} are the
standard references for the latter. Many applied studies, including the
benchmark instance used here, instead apply a single ``terrain occlusion
loss'' constant \cite{matolak,alhourani2014}. Section~\ref{sec:channel}
quantifies the bias this introduces into a relay-sizing decision: the demand side
rises, the supply side of admissible hover points falls, and the two compound.

\paragraph{Joint communication-and-trajectory design.}
Co-design of trajectories and communication is well studied for mobile relays
\cite{zeng2016relay,relay5g} and for multi-UAV and energy-harvesting systems
\cite{wu2018,you2019,xu2018wpt,zengtraj,mozaffari}. In those formulations the
degrees of freedom are the trajectory and the transmit power, and the objective is
a rate, an outage probability or an energy bill. The coupling here is different in
kind: the relay is quasi-static (hovering) and the transport aircraft is the
mobile terminal, so the interaction acts through the \emph{timing} of outages
rather than through throughput or interference. That is why our central object is
a blackout \emph{window} and a span function rather than a rate region, and why
the resulting feasibility condition is combinatorial --- a covering condition that
must hold in every demand cell --- rather than continuous.

\paragraph{Set covering, partitioning and the price of autonomy.}
The relay side is a covering problem in which the covering sets are no longer
static, which places it between classical set covering \cite{korte} and
continuous-time scheduling. The partitioning question we quantify in
Section~\ref{sec:partition} is the familiar observation that splitting a
resource pool duplicates its peak, here expressed as a weighted
demand ratio.

\paragraph{Position of this work.}
We retain the transport and relay layers of the applied literature but change
three things that alter the answers: an explicit handover model with a
checkable obstruction certificate; a physics-based diffraction model in place
of a constant penalty; and a feedback loop that re-optimises the transport
schedule for relay feasibility. Table~\ref{tab:ablation} collects the resulting
baselines and ablations.

%==============================================================================
\section{System model}
\label{sec:model}

\subsection{Instance and unified physical conventions}
A single depot $O$ serves $n_A=\ENnArea$ demand areas with $\ENnBox$ boxes.
The transport fleet has $\ENnDrone$ multirotors in three classes; the relay
fleet has $\ENnRelayFleet$ vehicles with $\ENnRelayBat$ energy packs. Links use
$f=\ENfMHz$\,MHz ($\lambda=\ENlambda$\,m) and the gateway antenna sits
$\ENgwAntH$\,m above ground.

Figure~\ref{fig:env} fixes the geography that drives everything downstream. The
relief is a Copernicus GLO-30 digital elevation model at $30$\,m posting
\cite{copernicusdem},
resampled to a $1/3600^\circ$ grid ($\ENnDemNx\times\ENnDemNy$ cells, elevations
$\ENdemMin$--$\ENdemMax$\,m); the depot sits in a valley at the south-centre and
the $\ENnArea$ demand areas fan out into terrain that rises by more than a
kilometre. That single fact --- a low depot and a high, dissected catchment ---
is what makes line-of-sight the binding constraint rather than payload or
endurance, and it is why Section~\ref{sec:channel} insists on a profile-based
diffraction model instead of a constant occlusion penalty.

\Dfig[0.90]{fig_en_env}{The instance. Thirty-metre Copernicus GLO-30
hillshade with $80$\,m contour lines, the depot O01 (star, valley floor) and the
$\ENnArea$ demand areas S001--S015 (dots); elevation scale on the right. The
relief that isolates the settlements is the same relief that occludes the
air-to-ground links: over the whole elevation model the ground spans
$\ENdemMin$--$\ENdemMax$\,m ($\ENdemRelief$\,m of relief) across a
$\sim\ENdemSpanKm$\,km-wide box, and inside the plotted window the terrain still
rises from $\ENdemWinMin$ to $\ENdemWinMax$\,m. The depot's valley position is
what makes the $\ENqCand$-point hover grid a genuinely constrained resource
rather than a formality.}{fig:env}

\subsection{Hovering-relay handover}
Let $\mathcal{C}$ be a finite set of admissible hover points (position plus
height above ground). A relay at $c\in\mathcal{C}$ covers a transport aircraft
at time $t$ if both the access link (aircraft$\to$relay) and the backhaul link
(relay$\to$gateway) close. We consider two relocation modes.

\begin{definition}[Blackout window]\label{def:blackout}
For $a,b\in\mathcal{C}$, $a\neq b$, the \emph{blackout window} $W(a,b)$ is the
time during which a single relay is in transit from $a$ to $b$ and therefore
covers nothing:
\begin{align}
W_{\mathrm{M1}}(a,b) &= t^{\mathrm{back}}(a) + \tau_{\mathrm{turn}}
  + t^{\mathrm{out}}(b) + t_{\mathrm{link}},\\
W_{\mathrm{M2}}(a,b) &= t^{\mathrm{transit}}(a\to b) + t_{\mathrm{link}},
\end{align}
where M1 returns to the depot to swap the energy pack and M2 transits directly
between hover points. $t^{\mathrm{out}},t^{\mathrm{transit}}$ follow the
climb--cruise--descend decomposition with cruise altitude set
$50$\,m above the highest terrain along the leg.
\end{definition}

\begin{proposition}[Elementary properties of $W$]
\label{prop:W}
If $t^{\mathrm{transit}}$ is symmetric, then for mode M2: (i) $W$ is symmetric,
$W(a,b)=W(b,a)$; (ii) $W(a,a)=t_{\mathrm{link}}$; and (iii) $W$ satisfies the
triangle-like bound $W(a,c)\le W(a,b)+W(b,c)-t_{\mathrm{link}}$. For mode M1,
$W$ is generally asymmetric and $W(a,a)=t^{\mathrm{back}}(a)+\tau_{\mathrm{turn}}
+t^{\mathrm{out}}(a)+t_{\mathrm{link}}>0$: \emph{even returning to the same
hover point costs a blackout}.
\end{proposition}
\noindent On our instance the empirical M1 blackout distribution over
$\ENqBlackoutPairs$ ordered pairs has median $\ENqBlackoutMedian$\,s
(min $\ENqBlackoutMin$\,s, max $\ENqBlackoutMax$\,s).

\subsection{Discrete demand}
Transport sorties are sampled at $10$\,s; an \emph{outage incidence} is a
sampled instant at which a sortie has no direct link to the gateway. Grouping
incidences by $10$\,s cell yields $\ENqInstances$ incidences on
$\ENqCells$ demand cells; up to $\ENqPeak$ aircraft may be disconnected
simultaneously.

\begin{definition}[Span function]\label{def:span}
For a demand profile $\mathcal{D}(\cdot)$, let
$\mathrm{cov}(\mathcal{D},c)\in\{0,1\}$ indicate that point $c$ alone covers all
incidences of cell $k$. The \emph{span} at cell $k$ is the largest number of
consecutive cells that some single point covers:
\eqfit{
\mathrm{span}(k)=\max\{\ell:\ \exists c\in\mathcal{C}\ \text{with }
\mathrm{cov}(\mathcal{D}(k'),c)=1\ \forall k'\in[k,k+\ell-1]\}.
}
We write $T$ for the length of the demand period, i.e.\ the elapsed time between
the first and the last demand cell. Two units appear in this paper and are kept
distinct: $\mathrm{span}(k)$ as defined above counts \emph{demand cells}, while
$\mathrm{span}_{\mathrm{s}}(k)$ denotes the elapsed seconds from the first to the
last instant of the same run. They are not proportional, because the cells are not
equally spaced (mean spacing $\approx 11.4$\,s), so every claim below states which
one it uses.
\end{definition}

\Dfigwide[0.96]{fig_en_q3_comm}{Link state of every transport sortie over the mission, from the shipped communication log ($236$ segments). Blue is a direct link to the gateway, green is a link carried by a hovering relay, red is an interval with no coverage at all. Two features drive the whole paper: outages are concentrated in long contiguous stretches rather than scattered samples, and the sorties overlap in time, so several of them can be uncovered at the same instant.}{fig:q3comm}

\subsection{Relay scheduling problem}
Given the demand profile, choose the number of relays $N$, their hover points
and the instants at which they relocate, so that every incidence is covered;
each relay must return before its energy pack is exhausted, and at most one
relay relocates at a time.

%==============================================================================
\section{Methodology}
\label{sec:method}

\subsection{Handover-aware feasibility}
\label{sec:handover}

\begin{theorem}[Handover obstruction for two relays]
\label{thm:handover}
Assume (A0) no \emph{configuration} covers the whole demand profile, where a
configuration assigns one hover point to each relay --- that is, neither a single
point nor any \emph{pair} of points covers the profile, so every feasible
two-relay schedule must relocate at least once; (A1) a relocating relay covers
nothing during its blackout window; and (A2) at most one relay relocates at a
time. Then a two-relay schedule is feasible \emph{only if} there exists a cell
$k$ with
\begin{equation}
\mathrm{span}\bigl(j_k\bigr)\ \ge\ k-j_k+1,
\qquad j_k = \min\{j:\ t_j \ge t_k-W(a,b)\},
\label{eq:star}
\end{equation}
for the pair $(a,b)$ of some handover, where $\mathrm{span}(\cdot)$ is the span of
Definition~\ref{def:span}, i.e.\ the number of
consecutive demand cells starting at $j$ that the stationary relay's point covers
alone. In the seconds convention this is simply
$W(a,b)\le\mathrm{span}_{\min}(k)$, with $\mathrm{span}_{\min}$ as defined below;
both forms are stated because the demand cells are \emph{not} equally spaced
(Assumption~\ref{asm:discrete}), so a cell count must not be converted to seconds
at a nominal $10$\,s rate. Consequently, if \eqref{eq:star} fails for \emph{every}
cell and every pair, no two-relay schedule exists.
\end{theorem}
\begin{proof}
Suppose a feasible two-relay schedule exists. By (A0) it must relocate at least
once, say the first relocation is from $a$ to $b$ and completes at instant $t$,
with the other relay stationary at $c$ throughout. By (A1) the moving relay
covers nothing on $[t-W,t]$, so by (A2) $c$ is the sole provider on that window.
Hence $c$ alone covers every demand cell whose instant lies in
$[t-W,t]$; writing $k$ for the last such cell, those cells are exactly
$j_k,\dots,k$ with $j_k=\min\{j: t_j\ge t_k-W\}$, so $\mathrm{span}(j_k)\ge
k-j_k+1$, which is \eqref{eq:star}. Note that (A0) is needed: if some pair
$\{a,b\}$ covered the whole profile, both relays could hover without ever
relocating and the conclusion would assert the existence of a handover that never
occurs.
\end{proof}

\noindent\textbf{Verifying (A0) on our instance.} A0 is an instance-level
hypothesis, not a modelling assumption, so we check it: no single candidate point
covers the whole profile (the longest solo span is $3970$\,s against a mission
length of $11290$\,s), and no \emph{pair} of points in $\mathcal{C}$ covers it
either, since the same consolidation test that drives
Proposition~\ref{prop:phase} finds no configuration whose union coverage is the
full profile. Both facts are recomputed by \texttt{code/span\_util.py}.

\noindent\textbf{What \eqref{eq:star} does and does not settle.}
Theorem~\ref{thm:handover} is a \emph{necessary} condition, and on our instance it
is far from binding: evaluating it against the actual pair $(a,b)$ of each
handover rather than against a fixed window, it is satisfiable at
$\ENqStarBase\%$ of the demand cells at the median M1 window and
$\ENqStarAir\%$ at the median M2 window, where the window lengths are the medians
of the full ordered-pair blackout distribution ($\ENqBlackoutPairs$ pairs) and the
satisfiability test itself is evaluated on a sampled subset of candidate pairs. We verified this in a
parameter-free form as well: writing
$\mathrm{span}_{\min}(k)=\max_{c}\{\ell: c \text{ covers } D \text{ on }
[t_k-\ell,t_k)\}$ for the longest window that \emph{some single point} can hold
alone at cell $k$, the condition reads $W(a,b)\le\mathrm{span}_{\min}(k)$ with no
tuned threshold; it holds for at least one pair in $880$ of the $992$ cells under
M1 and in $986$ of $992$ cells under M2. \emph{Therefore \eqref{eq:star} cannot
certify two-relay infeasibility, and we do not use it to do so}: the negative
answers reported below come from explicit search, and are stated with their search
scope. Conversely, the failure of \eqref{eq:star} at a specific cell is a
certificate, and the deficiency functional of
Corollary~\ref{cor:deficiency} measures how far the instance is from that
certificate's reach.

\begin{corollary}[Deficiency functional]
\label{cor:deficiency}
Work in the cell convention, so that indices are demand cells: let
$j_k=\max\{j:t_j\le t_k-W(a_k,b_k)\}$ be the first cell of the blackout window
charged to cell $k$, where $(a_k,b_k)$ is the handover pair charged to that cell.
Define the deficiency as the number of demand cells by which the window exceeds
what a single point can hold,
\eqfit{
D=\sum_{k}\max\bigl(0,\ (k-j_k+1)-\mathrm{span}(j_k)\bigr),
}
where $\mathrm{span}$ is the span of Definition~\ref{def:span}, i.e.\ the number
of consecutive cells that \emph{one} point covers. (Written in seconds the same
condition is $W(a_k,b_k)\le\mathrm{span}_{\min}(k)$ with
$\mathrm{span}_{\min}(k)=\max_{c}\{\ell:\ c \text{ covers } D \text{ on }
[t_k-\ell,t_k)\}$; we report $D$ in cells and do not convert it, for the reason
given below.) If hypotheses (A0)--(A2) of
Theorem~\ref{thm:handover} hold and \eqref{eq:star} fails everywhere, then
$D>0$ \emph{and} no two-relay schedule exists. $D$ is therefore an admissible
\emph{obstruction measure}: driving $D$ to zero is necessary for feasibility.
\end{corollary}

\noindent On our instance \eqref{eq:star} \emph{does} hold for
$\ENqStarBase\%$ of the cells under mode M1 ($\ENqStarAir\%$ under M2),
so the certificate does not fire: two-relay infeasibility is a genuinely
combinatorial obstruction rather than a scalar one. The deficiency $D$ of
Corollary~\ref{cor:deficiency}, summed over all cells, is $\ENqDeficitCellsBase$
demand cells under M1 and $\ENqDeficitCellsAir$ under M2. It is reported in demand
cells and never converted to seconds: the cells are not uniformly
spaced (mean spacing $\approx 11.4$\,s, Assumption~\ref{asm:discrete}), so a cell
count times a nominal $10$\,s would understate the elapsed time. Where an elapsed
duration is what matters, we state the seconds directly from the timestamps. For
$k$ whose window reaches before the mission starts the corresponding term
contributes $0$, since all demand in that range is zero.

\subsection{Exact feasibility characterisation}
\label{sec:theory}

Throughout we work under the following assumptions, which make the physical
layer a finite object without changing the logic of the arguments.

\begin{assumption}[Handover cost]\label{asm:handover}
A relay that relocates from $a$ to $b$ provides no service during the blackout
window $W(a,b)$ of Definition~\ref{def:blackout}. Relocation is otherwise instantaneous at its
completion instant.
\end{assumption}
\begin{assumption}[Serialised relocation]\label{asm:serial}
At most one relay relocates at any instant; the others hold their hover points.
\end{assumption}
\begin{assumption}[Energy]\label{asm:energy}
A relay may remain airborne for at most $T^{\mathrm{cap}}$ seconds between two
visits to the depot, and must return before its state of charge falls below the
reserve fraction $\rho$.
\end{assumption}
\begin{assumption}[Recharge gap]\label{asm:recharge}
A relay that has returned to the depot to swap energy modules provides no service
until it is airborne again; the depot offers no alternative coverage.
\end{assumption}
\begin{assumption}[Discretisation]\label{asm:discrete}
The demand profile is evaluated on a finite increasing sequence of
\emph{demand cells} $t_0<\dots<t_{n-1}$; cells need not be equally spaced.
\end{assumption}

\begin{theorem}[Exact reachability]
\label{thm:exact}
Disable all pruning. The state recursion of Section~\ref{sec:handover} reaches a
non-empty set of valid states at the last demand cell if and only if an
$N$-relay schedule satisfying Assumptions~\ref{asm:handover}--\ref{asm:discrete} exists.
\end{theorem}
\begin{proof}
($\Rightarrow$) Any reachable state path is, by construction, a schedule: each
transition relocates one relay subject to the blackout, coverage and energy
conditions, and validity at every cell means every incidence is covered.
($\Leftarrow$) Given a feasible schedule, order its relocation completions by
time; each corresponds to one transition, and the intermediate configurations
are exactly the states visited. All conditions (V), (T1)--(T5) hold at each
step, so the recursion cannot miss the path.
\end{proof}

Algorithm~\ref{alg:dp} states the recursion exactly as it is implemented.
Its only approximations are the ones discussed after Theorem~\ref{thm:exact}:
a finite candidate set and, in the runnable version, state pruning.

\begin{algorithm}[htbp]
\caption{Handover-aware relay feasibility (stated without pruning; the
runnable implementation adds the pruning reported in
Section~\ref{sec:reproducible})}
\label{alg:dp}
\begin{algorithmic}[1]
\Require demand cells $t_0<\dots<t_{n-1}$ with sets $\mathcal{D}_k$; candidate
points $\mathcal{C}$ with coverage $\mathrm{cov}[k,c]$; blackout $W$; cap
$T^{\mathrm{cap}}$; fleet size $N$
\Ensure a feasible schedule, or the first cell with no feasible state
\State $S_0 \gets \bigl\{\,((\mathrm{IDLE},t_0,t_0))_{i=1}^{N}\,\bigr\}$
\For{$k=0$ to $n-1$}
  \State $S_k \gets \{\,s\in S_k:\ \bigcup_{i:\,p_i\neq\mathrm{IDLE}}\mathrm{cov}[\cdot,p_i]\supseteq\mathcal{D}_k\,\}$ \Comment{(V)}
  \If{$S_k=\emptyset$} \State \Return ``no feasible state at cell $k$''
  \EndIf
  \ForAll{$s\in S_k$, $i\in\{1..N\}$, $p'\in\mathcal{C}\cup\{\mathrm{IDLE}\}$, $p'\neq p_i$}
    \State $j \gets j_k \gets \textsc{Before}(k,\,W(p_i,p'))$ \Comment{last cell with $t_{j}\le t_k-W$; not $\lceil W/\Delta\rceil$}
    \If{(T1) $\sigma_i\le t_{j}$ \textbf{and} (T2) $\sigma_j\le t_{j}\ \forall j\neq i$}
      \If{(T3) $\bigcup_{j\neq i}\mathrm{cov}[\cdot,p_j]\supseteq\mathcal{D}_{k'}\ \forall k'\in[j+1,k-1]$}
        \If{(T4) $\bigcup_{j}\mathrm{cov}[\cdot,p_j]\supseteq\mathcal{D}_k$ \textbf{and}
             (T5) $t_k-\varrho_i^{(k)}\le T^{\mathrm{cap}}\ \forall i$}
          \State $S_{k+1}\gets S_{k+1}\cup\{\,s\ \text{with}\ (p_i,\sigma_i)\gets(p',t_k),\ \varrho_i\gets t_k\ \text{iff mode M1 or }p_i=\mathrm{IDLE}\,\}$
        \EndIf
      \EndIf
    \EndIf
  \EndFor
\EndFor
\State \Return any state path in $S_{n-1}$
\end{algorithmic}
\end{algorithm}

\noindent\textbf{Cost.} One transition sweep touches $|S_k|\cdot N\cdot|\mathcal{C}|$
candidates and evaluates (T3) over the $k-j_k$ cells of the blackout window, so
the per-cell work is $O(|S_k|\,N\,|\mathcal{C}|\,(k-j_k))$; without pruning
$|S_k|$ grows as $|\mathcal{C}|^{N}$. This is why the runnable implementation
prunes, and why the pruning policy is reported with every infeasibility claim.

{\sloppy\noindent\textbf{Why the window is indexed by timestamps.} The demand cells are
\emph{not} a uniform $10$\,s grid: they record only the instants at which some
transporter is out of direct range, and their mean spacing is about $11.4$\,s
(Assumption~\ref{asm:discrete}). Converting a blackout window $W$ to a nominal
cell count $\lceil W/\Delta\rceil$ and then counting cells would therefore
measure the window as roughly $\lceil W/10\rceil\times 11.4$\,s --- an error of
about $14\%$ that grows with the window.}\par Algorithm~\ref{alg:dp} therefore
locates the window by the timestamp search $$\textsc{Before}(k,W')=\max\{j:
t_j\le t_k-W'\},$$ which is what the implementation does, and the seconds form
$W\le\mathrm{span}_{\min}(k)$ of \eqref{eq:star} is the same condition written
without reference to cells at all.

\noindent Theorem~\ref{thm:exact} is the reason we can speak of a
\emph{decision} rather than a search failure: any implementation that prunes
states forfeits the ``only if'' direction, which is why all infeasibility claims
in this paper are stated as ``no feasible schedule found under the stated
candidate set and pruning policy''.

\begin{corollary}[Continuous-time version]
\label{cor:continuous}
Replace the demand cells by a measurable demand set $D(t)$ and let
$\mathrm{span}(\tau)=\sup\{\ell:\exists c\in\mathcal{C},\ c\ \text{covers}\
D(t)\ \forall t\in[\tau,\tau+\ell)\}$. If no single point covers
$D(t)$ for all $t$, then a two-relay schedule under
Assumptions~\ref{asm:handover}--\ref{asm:serial} exists only if there is a handover instant $\tau$ and a pair
$(a,b)$ with $\mathrm{span}(\tau-W(a,b))\ge W(a,b)$.
\end{corollary}
\begin{proof}
Identical to Theorem~\ref{thm:handover}; discreteness is used nowhere in that
argument.
\end{proof}

\begin{theorem}[Exact decision for two relays]
\label{thm:two}
Under Assumptions~\ref{asm:handover}--\ref{asm:discrete}, fix a candidate set
$\mathcal{C}_0\subseteq\mathcal{C}$. A two-relay schedule with hover points in
$\mathcal{C}_0$ exists \emph{if and only if} there is a sequence of
configurations $(a_0,b_0)\to(a_1,b_1)\to\cdots\to(a_m,b_m)$ with
$a_r,b_r\in\mathcal{C}_0$ and completion instants
$0\le\tau_1<\cdots<\tau_m\le T$ such that, writing $W_r=W(a_{r-1},a_r)$ and
$W'_r=W(b_{r-1},b_r)$ for the blackout windows of the moves:
\begin{enumerate}[leftmargin=1.8em,itemsep=2pt]
\item[(C1)] at most one relay is in blackout at any instant:
  Assumption~\ref{asm:serial} forbids simultaneous relocation, so if the two
  windows overlap we may delay the later completion until the earlier window has
  ended without changing the demand seen by either relay. Moves need \emph{not}
  alternate between the relays --- one relay may complete several relocations in
  succession while the other holds its point, and the recursion of
  Section~\ref{sec:handover} searches that case as well;
\item[(C2)] for every $r$, the \emph{stationary} relay covers the whole demand
  over the blackout window of the moving one: $\mathrm{cov}[\cdot,b_{r-1}]$
  alone covers $D(t)$ for all $t\in[\tau_r-W_r,\tau_r]$ when the first relay
  moves, and symmetrically for the second;
\item[(C3)] between two consecutive completions of the same relay the
  configuration is constant, and the sortie length is within $T^{\mathrm{cap}}$.
\end{enumerate}
Consequently, on a fixed $\mathcal{C}_0$, the two-relay question is decided by
the state recursion of Theorem~\ref{thm:exact} with pruning disabled: it returns
either a schedule or an empty state set, and the latter is a proof of
infeasibility over $\mathcal{C}_0$.
\end{theorem}
\begin{proof}
($\Rightarrow$) Take a feasible schedule and list its relocation completions in
time order. By Assumption~\ref{asm:serial} at most one relay moves at a time, so
the relocation windows are pairwise disjoint, which is (C1) --- note that (C1)
does \emph{not} require the relays to take turns: the same relay may relocate
twice in a row, and the construction below allows it. (C2) holds
because the moving relay contributes nothing (Assumption~\ref{asm:handover})
while the stationary one keeps its point; (C3) restates the energy rule.
($\Leftarrow$) Given such a sequence, hold each configuration between
completions. Conditions (C1)--(C2) make every interval covered---during a window
the stationary relay covers by (C2), and outside windows the standing
configuration covers by construction---while (C3) makes the energy feasible, so
the schedule is admissible. Finally, with pruning disabled the recursion of
Theorem~\ref{thm:exact} enumerates exactly the sequences described by
(C1)--(C3), which gives the decision procedure.
\end{proof}

\noindent The practical value of Theorem~\ref{thm:two} is that the negative half of
the fleet-size question can be turned from an open-ended search failure into a list
of blocked handovers. Two consequences are used below. First, the certificate is
\emph{decidable on a restricted candidate set}: every infeasibility statement we
report about $N=2$ is made over an explicitly listed candidate set and a stated
pruning policy, so it is a statement about a finite search space rather than about
search effort (Proposition~\ref{prop:distinguishing}). We stress that the shipped
run is the \emph{pruned} one: closing state pruning at this breadth did not
terminate within our budget, so the verdict is relative to the declared set and we
do not upgrade it to an unpruned decision. Second, the two-relay case
is special only in cost, not in logic: for general $N$ the same recursion decides
the question, but its state set grows like $|\mathcal{C}|^{N}$
(Section~\ref{sec:theory}), which is why the constructive verifier of
Section~\ref{sec:verifier} exists. We therefore always report the candidate set
\emph{and} the per-cell candidate refinement together with an infeasibility
claim, since the refinement is the one layer that Proposition~\ref{prop:distinguishing}
does not cover.

\begin{proposition}[Distinguishing candidate set]
\label{prop:distinguishing}
Fix a candidate set $\mathcal{C}_0\subseteq\mathcal{C}$ and run the recursion of
Theorem~\ref{thm:exact} on $\mathcal{C}_0$ with pruning disabled. If it generates
no admissible state at some cell, then no two-relay schedule exists whose hover
points are drawn from $\mathcal{C}_0$. This conclusion is unconditional in
$\mathcal{C}_0$: it is a statement about a finite, explicitly listed search space
and does not depend on the enumeration order, on any pruning rule, or on the
number of states retained.
\end{proposition}
\begin{proof}
Let $G_{\mathcal{C}_0}$ be the graph whose vertices are the admissible states and
whose edges are the transitions of the recursion, both restricted to
$\mathcal{C}_0$. By construction every edge corresponds to a physically
admissible relocation and every vertex to a configuration that covers its cell;
conversely, any schedule with points in $\mathcal{C}_0$ yields, after ordering its
relocation completions, a walk in $G_{\mathcal{C}_0}$ that visits one vertex per
demand cell. The recursion with pruning disabled performs an exhaustive
expansion of $G_{\mathcal{C}_0}$, so an empty state set at cell $k$ means that no
walk from the initial vertex reaches cell $k$, hence no such schedule exists.
\end{proof}

\begin{proposition}[Energy lower bound]
\label{prop:energy}
Assume the coverage requirement is \emph{continuous}: every instant of
$[t_0,t_{n-1}]$ must be covered by some airborne relay, and by
Assumption~\ref{asm:recharge} a relay at the depot covers nothing. Then
$N\ge\bigl\lceil T/T^{\mathrm{cap}}\bigr\rceil$ with $T=t_{n-1}-t_0$.
\end{proposition}
\begin{proof}
Let $A_i\subseteq[t_0,t_{n-1}]$ be the set of instants at which relay $i$ is
airborne and therefore able to provide coverage. Continuous coverage means
$[t_0,t_{n-1}]=\bigcup_{i=1}^{N}A_i$, so $T\le\sum_i|A_i|$. By
Assumption~\ref{asm:energy} each maximal interval of $A_i$ has length at most
$T^{\mathrm{cap}}$, but $A_i$ may consist of $m_i>1$ such intervals, so
$|A_i|\le m_iT^{\mathrm{cap}}$ and the decomposition alone yields only
$T\le T^{\mathrm{cap}}\sum_i m_i$, which says nothing about $N$: a single relay
recharging repeatedly is not excluded by Assumption~\ref{asm:energy}. The bound
therefore requires one further hypothesis, which we state explicitly and verify on
our instance rather than leave implicit:

\smallskip
\noindent\emph{(H) the coverage intervals of distinct relays can be taken
pairwise disjoint.}
\smallskip

\noindent Under (H) at most one relay is covering at any instant, the airtime of
distinct relays does not overlap in the covering role, and
$T=\sum_i|A_i|\le\sum_iT^{\mathrm{cap}}=N\,T^{\mathrm{cap}}$, giving
$N\ge\lceil T/T^{\mathrm{cap}}\rceil$.

\smallskip
\noindent On our instance (H) is exactly what the mission demands rather than a
convenience: coverage must be sustained at $\ENqCells$ demand cells spanning
$T=\ENTs$\,s, and the verified schedule at $N=\ENrbGreedy$ uses
$\ENgrNThreeMOneSorties$ sorties with SOC minimum $\ENgrNThreeMOneSoc\%$, i.e.\ the
relays \emph{hand the covering role over} instead of overlapping it. We therefore
report $N\ge\ENboundEnergy$ as a bound \emph{conditional on (H)}, and we do not
rely on it for any fleet-size claim that the constructive schedule does not
already support.
\end{proof}

\noindent The hypotheses of both propositions are verified on our instance by
direct computation. The common premise (A0) of Theorem~\ref{thm:handover} is
that no configuration covers the whole profile. For a single point this is
immediate: the longest stretch that any one hover point can hold alone is
$\ENspanMax$\,s against a mission length of $\ENTs$\,s, a factor of
$\ENTs/\ENspanMax$ short. For a \emph{pair} we test the consolidated coverage of
every candidate pair over the full profile, and no pair covers it either --- the
same test that drives Proposition~\ref{prop:phase}. For
Proposition~\ref{prop:energy} the additional hypothesis (H) is met by the verified
schedule (the relays hand the covering role over rather than overlap it), and
$T^{\mathrm{cap}}=\ENcapS$\,s then yields $N\ge\ENboundEnergy$. We stress that this
is a \emph{conditional} bound and that the two-relay question is \emph{not} closed
by it --- for that we rely on the search reported in
Section~\ref{sec:premium}, with its scope stated.

\begin{proposition}[Relocation counting bound]
\label{prop:structural}
Let $w_{\min}=\min_{a\neq b}W(a,b)>0$ and let $r$ be the total number of
relocation completions in a feasible schedule. Then
$r\le 1+\lfloor T/w_{\min}\rfloor$. Moreover, if no \emph{configuration} (a
choice of one hover point per relay) covers the whole demand profile, then every
relay must relocate, so $N\le 1+\lfloor T/w_{\min}\rfloor$ is required for
feasibility to be possible at all.
\end{proposition}
\begin{proof}
Assumption~\ref{asm:serial} makes the blackout windows of all relocations
pairwise disjoint; each has length at least $w_{\min}$ and each completion lies in
$[0,T]$, so the completions are separated by at least $w_{\min}$ and
$r\,w_{\min}\le T$. It remains to connect $N$ to $r$. By the premise, no single
configuration (one hover point per relay) covers the whole profile, so for each
relay $i$ there is a demand cell that $i$'s point fails to cover; covering that
cell therefore requires some relay to be elsewhere, i.e.\ to relocate. This holds
for every relay, so all $N$ relocate and $r\ge N$. Combining gives
$N\le r\le 1+\lfloor T/w_{\min}\rfloor$.
\end{proof}

\noindent The bound is a screening test rather than a sizing rule: on the instance
it reads $N\le\ENboundCounting$ with $w_{\min}=\ENwMinMOne$\,s under M1, which is
far from binding. It is reported because it explains why the fleet size is
\emph{not} limited by relocation throughput here, and therefore why the answers
must come from a state-space argument (Theorem~\ref{thm:two}) rather than from a
counting one.

\begin{proposition}[Phase-level relaxation is not conservative]
\label{prop:phase}
Partition the demand period into phases and consider the relaxation in which a
relay may choose one hover point \emph{per phase} (rather than per instant),
subject to (i) every incidence being covered by some relay's phase point and
(ii) $W(a,b)$ fitting into the interval between the end of phase $k$ and the
start of phase $k+1$ whenever a relay uses points $a,b$ in those two phases.
Call a schedule \emph{phase-stationary} if every relay keeps one hover point
throughout each phase. If a phase-stationary schedule exists for $N$ relays, the
relaxation is feasible for $N$ relays. Two caveats follow, and both matter:
(a) the hypothesis is not automatic --- a schedule may relocate \emph{inside} a
phase, in which case the induced per-phase assignment is bounded by
$2N$ points rather than $N$ and the implication is not immediate; and
(b) the converse fails in general, and it fails on our instance: the relaxation is
feasible for $N=2$ (witnessed by an explicit assignment, verified by exhaustive
search over the relaxed model) while the exact recursion of
Theorem~\ref{thm:two} returns no feasible state.
\end{proposition}
\begin{proof}
{\sloppy\emph{Forward.} Given a phase-stationary schedule, assign to each relay
and each phase the point at which that relay hovers during that phase.}\par A relay serves at
most one point per phase by hypothesis, and every incidence is covered because the
schedule covers every instant. It remains to check (ii). If a relay uses different
points in two consecutive phases, it must relocate; by
Assumption~\ref{asm:handover} the relocation occupies an interval of length at
least $W(a,b)$ during which the relay covers nothing, and by
Assumption~\ref{asm:serial} that interval may not overlap another relay's
relocation. Since the relay must be in place for the whole of phase $k+1$, the
interval must finish by the start of phase $k+1$; since it must have left its
phase-$k$ point, it starts no earlier than the end of phase $k$ at which it last
held that point. Hence $W(a,b)$ is at most the length of the inter-phase gap, which
is condition (ii). Hence the relaxed model admits the induced assignment.
{\sloppy\emph{Converse.} The relaxation drops the requirement that the
\emph{stationary} relay hold a single point throughout the moving relay's
blackout window; on our instance this extra requirement is what makes $N=2$
infeasible while the relaxation is satisfiable.}\par
\end{proof}

\noindent We use this proposition only in the direction that survives the caveat
above: a \emph{negative} answer from the relaxation is informative (it rules out
phase-stationary schedules of that size), while a \emph{positive} answer is weak
because it may exploit a relocation that the exact model forbids. That is exactly
why we report the phase-level integer program \emph{alongside} rather than
\emph{instead of} the exact recursion.

\noindent The single-relay case is the boundary of the problem and is worth
stating exactly, because it is the only fleet size at which the answer is a
closed-form condition rather than a search: with $N=1$ no relocation is possible
at all, so feasibility collapses to a property of one hover point. The theorem
below is the $N=1$ specialisation that the recursion of
Theorem~\ref{thm:two} reproduces, and it is what makes the $N=1$ row of
Table~\ref{tab:ilp} an independent check rather than a repetition.

\begin{theorem}[Single relay]
\label{thm:one}
Assume a single relay and let $t_{\mathrm{f}}$ (resp.\ $t_{\ell}$) be the first
(resp.\ last) instant of demand. A feasible schedule exists if and only if there
is a hover point $c\in\mathcal{C}$ such that
\begin{enumerate}[leftmargin=1.6em,itemsep=2pt]
\item $c$ covers $D(t)$ for every demand instant $t\in[t_{\mathrm{f}},t_{\ell}]$, and
\item $t^{\mathrm{out}}(c)+(t_{\ell}-t_{\mathrm{f}})+t^{\mathrm{back}}(c)
 \le T^{\mathrm{cap}}$.
\end{enumerate}
\end{theorem}
\begin{proof}
A single relay cannot relocate without leaving the demand uncovered, because
relocation always costs a strictly positive blackout. That positivity is part
(ii) for M2 ($W(a,a)=t_{\mathrm{link}}>0$, so even a no-op relocation blacks the
link out) and the closing sentence of Proposition~\ref{prop:W} for M1
($W(a,a)>0$ there as well); note that it is \emph{not} part (iii), which is the
triangle bound and says nothing about sign. Hence its hover
point must be constant throughout $[t_{\mathrm{f}},t_{\ell}]$, which is exactly
condition 1, and the sortie duration is exactly the expression in condition 2.
\end{proof}

\begin{proposition}[Monotonicity in fleet size]
\label{prop:monotone}
Let $\mathcal{F}(N)$ denote feasibility with $N$ relays. Then
$\mathcal{F}(N)\Rightarrow\mathcal{F}(N+1)$.
\end{proposition}
\begin{proof}
Keep the schedule of the $N$ relays and let the additional relay remain at the
depot; it neither helps nor hinders.
\end{proof}

\begin{proposition}[Monotonicity in demand and in the energy budget]
\label{prop:demand}
Let $D\subseteq D'$ be two demand profiles on the same cell grid. Then
$\mathrm{span}'\le\mathrm{span}$ pointwise, the deficiency satisfies
$D'\ge D$, and $\mathcal{F}(N;D')\Rightarrow\mathcal{F}(N;D)$. Likewise,
enlarging $T^{\mathrm{cap}}$ or relaxing $\rho$ can only enlarge the set of
feasible schedules.
\end{proposition}
\begin{proof}
A point that covers $D'$ on an interval covers $D$ on it, so every span value
weakly decreases; the deficiency is a sum of pointwise non-increasing terms.
Any schedule that is valid for $D'$ is valid for $D$, since validity is a
covering condition. The energy statement is immediate from
Assumption~\ref{asm:energy}.
\end{proof}

\noindent Proposition~\ref{prop:demand} is what makes the terrain study of
Section~\ref{sec:exp} meaningful: the fleet size is monotone in demand
\emph{set inclusion}, yet the experiments show it is \emph{not} monotone in
scalar summaries of the demand such as the occlusion ratio. For the same reason
the recursion is run upwards from the smallest fleet that can conceivably work:
Proposition~\ref{prop:monotone} guarantees that if $N$ relays suffice then so do
$N+1$, so a negative verdict at some $N$ is not refuted by a positive verdict at a
larger one, and the reported $N^{*}$ is the first size that succeeds rather than a
global minimum over an unordered family.

\begin{proposition}[Regularity of the blackout function]
\label{prop:lipschitz}
Fix the hover heights. Let $M$ bound the terrain slope, $|\nabla z|\le M$, and
let $v_{\uparrow},v_{\downarrow},v_{c}$ be the climb, descent and cruise speeds.
Then $W$ is Lipschitz in the endpoints,
\eqfit{
|W(a,b)-W(a',b')|\ \le\ L\bigl(|a-a'|+|b-b'|\bigr),\qquad
L=\Bigl(\tfrac{1}{v_{\uparrow}}+\tfrac{1}{v_{\downarrow}}\Bigr)(1+M)+\tfrac{1}{v_{c}},
}
for mode M2, and the same bound augmented by the depot legs holds for M1.
\end{proposition}
\begin{proof}
The cruise altitude is $\max\{\text{line maximum},\ \text{endpoint altitudes}\}+50$,
a maximum of $1$-Lipschitz and $M$-Lipschitz functions of the endpoints, hence
$(1+M)$-Lipschitz; the climb and descent heights are this quantity minus a
constant, and the horizontal distance is $1$-Lipschitz by the triangle
inequality. Composing with the linear time map gives the constant $L$.
\end{proof}

\noindent Proposition~\ref{prop:lipschitz} has a practical reading: a small
positioning error of a hovering relay changes the blackout window by at most
$L$ seconds per metre, so the handover schedule is not brittle with respect to
hover-point placement---an observation that matters because the candidate set
$\mathcal{C}$ is a finite grid.

\subsection{Constructive verifier and coverage baseline}
\label{sec:verifier}
A feasible $N$ is \emph{proved} by exhibiting a full schedule; we therefore use a
constructive verifier (in the spirit of large-neighbourhood repair heuristics
\cite{alns}) that advances cell by cell and, whenever the current
configuration cannot serve the next cell, relocates one relay subject to the
blackout constraints. Two heuristic variants are run (with and without
\emph{pre-positioning} of an idle relay before the incumbent's energy deadline);
whichever yields a feasible schedule certifies feasibility.

Algorithm~\ref{alg:greedy} is the constructive verifier. Two heuristic
variants are run --- with and without the pre-positioning branch (lines
\ref{line:pre}--\ref{line:preend}) --- and whichever returns a complete schedule
certifies feasibility. The pre-positioning branch was added after diagnosing a
terrain realisation in which a single, long, \emph{sparse} outage exhausted the
incumbent relay's airborne budget: the successor must be dispatched late enough
to arrive fresh, but early enough to be in place, which is exactly the timing
window tested on lines \ref{line:pre}--\ref{line:preend}.

\begin{algorithm}[htbp]
\caption{Constructive verifier for $N$ hovering relays}
\label{alg:greedy}
\begin{algorithmic}[1]
\Require as in Algorithm~\ref{alg:dp}; per-cell candidate budget $b$; safety factor $\eta$
\Ensure a complete schedule (a proof of feasibility) or the first failing cell
\State $k\gets0$;\ \ $p_i\gets\mathrm{IDLE}$,\ $\mathrm{arr}_i\gets0$,\ $\varrho_i\gets0$ for all $i$
\While{$k<n$}
  \If{$\bigcup_i \mathrm{cov}[\cdot,p_i]\supseteq\mathcal{D}_k$ \textbf{and} $\varrho$ within $T^{\mathrm{cap}}$}
    \State $k_{\mathrm{dead}}\gets$ first cell at which some relay exceeds $T^{\mathrm{cap}}$
    \If{$k_{\mathrm{dead}}$ exists, an idle relay exists, and $k_{\mathrm{dead}}-k\le$ lead}
    \label{line:pre}
      \State $p'\gets\arg\max_{p\in\mathrm{cand}(k_{\mathrm{dead}})}\ \mathrm{run}(p,k_{\mathrm{dead}})$
      \State $t_{\mathrm{arr}}\gets t_k+t^{\mathrm{out}}(p')+t_{\mathrm{link}}$
      \If{$t_{\mathrm{arr}}\le t_{k_{\mathrm{dead}}}$ \textbf{and}
          $t_{k_{\mathrm{dead}}}-t_{\mathrm{arr}}\le \eta\,T^{\mathrm{cap}}$}
        \State dispatch the idle relay to $p'$; \ $k\gets k+1$; \textbf{continue}
      \EndIf
    \label{line:preend}
    \EndIf
    \State $k\gets k+1$; \textbf{continue}
  \EndIf
  \State {\sloppy enumerate admissible moves $(i,p')$ satisfying (T1)--(T5) and
         score them by $(-\text{urgency},\,-\text{run length},\,W)$, keeping the
         best one}\par
  \If{none is admissible} \State \Return $k$ \Comment{first failing cell}
  \EndIf
  \State perform the move; \ $k\gets k+1$
\EndWhile
\State \Return the reconstructed sortie list
\end{algorithmic}
\end{algorithm}

\noindent The non-executable baseline ignores handovers entirely: for each cell $k$ let
$c(k)$ be the minimum number of points whose union covers $\mathcal{D}(k)$
(exact set cover, a classical NP-hard problem \cite{korte}, solved here by
CP-SAT \cite{cpsat}), and set
$N_{\mathrm{naive}}=\max_k c(k)$. On our instance $N_{\mathrm{naive}}=\ENrbNaive$
while the executable model requires three relays for every schedule we succeeded
in constructing ($N^{*}\le\ENrbGreedy$, with $N^{*}\ge\ENboundEnergy$ proven): the handover premium is
$\ENrbPremium$ vehicles.

\subsection{Channel models}
\label{sec:channel}
\paragraph{Baseline (as supplied).}
A fixed penalty $L_{\mathrm{obs}}=\ENlObs$\,dB is added whenever the LoS is
intersected, and the link closes if $\mathrm{FSPL}+L_{\mathrm{obs}}>L_{\max}$.
Two structural defects follow: a short blocked link may still close
(e.g.\ at $2$\,km, $\mathrm{FSPL}\approx 106$\,dB), and a grazing ridge costs
the same as a mountain.

\paragraph{ITU-R P.~526 single knife edge.}
{\sloppy Along the terrain profile we compute, for every interior sample, the
Fresnel--Kirchhoff parameter
$\nu=h\sqrt{2(d_1+d_2)/(\lambda d_1 d_2)}$ \cite{iturp526} with $h$ the height of
the obstacle above the (Earth-curvature-corrected) LoS, $d_1,d_2$ the distances
to the endpoints, and take the dominant edge.}\par The excess loss is
\eqfit{
J(\nu)=6.9+20\lg\!\Bigl(\sqrt{(\nu-0.1)^2+1}+\nu-0.1\Bigr)\ \mathrm{dB},
\qquad \nu>-0.78,
}
and $J(\nu)=0$ otherwise \cite{iturp526}. The total loss is
$\mathrm{FSPL}+J(\nu_{\max})$, where FSPL is the free-space loss of
\cite{friis1946}. We also report a conservative multi-edge upper bound
$\sum_{\nu>0}J(\nu)$ after the Deygout construction \cite{deygout} and a $k=4/3$
effective-Earth-radius correction $d_1d_2/(2kR)$ \cite{iturp526}.

\subsection{Co-design feedback loop}
\label{sec:feedback}
Only the departure instants $\{s_i\}$ of the transport sorties are varied;
batching, routes, aircraft types and assignments are frozen. We minimise the
deficiency $D(s)$ of Corollary~\ref{cor:deficiency} by threshold accepting over
per-sortie feasible windows, and we evaluate the transport side
(makespan, weighted tardiness, hard-deadline violation) exactly at every
accepted move.

Algorithm~\ref{alg:code} minimises the deficiency of
Corollary~\ref{cor:deficiency} over departure instants only.

\begin{algorithm}[htbp]
\caption{Co-design feedback loop (threshold accepting on departure instants)}
\label{alg:code}
\begin{algorithmic}[1]
\Require frozen batching, routes, aircraft types and assignments; feasible
departure window $[l_i,u_i]$ and duration $d_i$ for each sortie $i$ (so the
latest departure that still finishes inside the window is $u_i-d_i$); deficiency
oracle $D(\cdot)$; step scale $\sigma$
\Ensure best departure vector and its deficiency
\State $s\gets s^{0}$;\ \ $D^{*}\gets D(s^{0})$;\ \ $s^{*}\gets s^{0}$
\For{$r=1$ to $R$} \Comment{restarts}
  \State $s\gets$ perturb($s^{0}$)
  \For{$\mathrm{it}=1$ to $I$}
    \State pick $i$; \ $s'\gets s$ with $s_i\gets\mathrm{clip}(s_i+\mathcal{N}(0,\sigma),\,l_i,\,u_i-d_i)$
    \State $T\gets T_0\,(1-\mathrm{it}/I)$
    \If{$D(s')\le D(s)$ \textbf{or} $\mathrm{rand}()<\exp\!\bigl(-(D(s')-D(s))/T\bigr)$}
      \State $s\gets s'$
    \EndIf
    \If{$D(s)<D^{*}$} \State $D^{*}\gets D(s)$;\ \ $s^{*}\gets s$ \EndIf
  \EndFor
\EndFor
\State \Return $s^{*},D^{*}$
\end{algorithmic}
\end{algorithm}

\noindent The deficiency oracle is evaluated by mapping every outage sample to
its demand cell, forming the (cell\,$\times$\,point) matrix of single-point
coverage, taking the longest run per point (\texttt{span\_util}\allowbreak
\texttt{.span\_cells}), and
then applying the time-window test of Corollary~\ref{cor:deficiency}. The
transport cost (makespan, weighted tardiness, hard-deadline violation) is
evaluated exactly at every accepted move, so no Pareto point in
Figure~\ref{fig:feedback} is estimated.

\noindent\textbf{Sample size and its consequence.} This experiment moves one
coordinate of the design space, and the reader should keep that limit in view.
The loop perturbs \emph{departure instants only}: batching, routes, aircraft
types and sortie-to-vehicle assignments are frozen at the baseline plan, so the
neighbourhood explored is the box $[l_i,u_i-d_i]$ around that plan. The search
uses $\ENcfRestarts$ restarts of $\ENcfIters$ threshold-accepting iterations
each, and the deficiency oracle is evaluated at discretised departure instants
sampled on the continuous box $[l_i,u_i-d_i]$ rather than on the continuous domain
itself. The reported response is therefore flat \emph{on that
neighbourhood, at that resolution}: the deficiency stays at
$\ENcfBaseDefA\to\ENcfBaseDefB$ demand cells under M1 and
$\ENcfAirDefA\to\ENcfAirDefB$ under M2, with $\ENcfBaseShiftN$ sorties actually
moved. A study that re-optimised routes and fleet composition jointly would range
over a strictly larger set and could in principle do better; we claim nothing
against that. What the null result does support is the narrower statement used in
the Discussion: on the frozen baseline plan, retiming departures alone does not
close the handover deficiency.

\subsection{Controlled terrain family}
\label{sec:terrain}
Because only one digital elevation model is available, we hold the node layout,
cargo and parameters fixed and transform the elevation field:
relief scaling $z' = z_{\mathrm{med}}+\alpha(z-z_{\mathrm{med}})$, corridor
elevation (a $1.8$\,km-wide corridor whose walls are raised by up to $180$\,m),
and \emph{terrain translation} $z''(i,j)=z'(i+\Delta i,j+\Delta j)$, which
changes which obstacle blocks which leg at constant relief. The occlusion ratio
-- the fraction of sampled transport-track instants without a direct link --
is used as the measurable abscissa.

%==============================================================================
\Dfigwide[0.96]{fig_en_q1_pareto}{The transport frontier and where the named plans sit. Filled markers are the Pareto frontier of the single-depot routing problem, coloured by sortie count; open circles are the eighteen named plans produced by the baselines and ablations of Section~\ref{sec:ablation}. Every named plan is dominated on both axes: the frontier trades energy against makespan, and the plans we compare against in the tables fall well inside it.}{fig:q1pareto}

\section{Experiments}
\label{sec:exp}

\subsection{Setup}
All experiments share one frozen transport plan and one candidate generation
rule, so that only the factor under study varies. The plan comprises
$\ENnTransport$ sorties carrying $\ENnBox$ boxes, with makespan
$\ENmakeSpan$\,s, transport energy $\ENtransEnergy$\,kWh and on-time rate
$\ENontimeRate\%$. Every result below is regenerated from the shipped result
files.

\Dfigwide[0.96]{fig_en_q2_gantt}{The frozen transport plan, sortie by sortie. Each bar is one sortie on one drone and they do not overlap, so the drone and battery resources are actually respected rather than assumed; the hatched tail of a bar is the battery turnaround that blocks the next sortie on that drone. Dashed vertical lines mark the hard deadlines of the medical and first-batch cargo. The relay layer of Sections~\ref{sec:premium}--\ref{sec:ilp} is evaluated against exactly this schedule.}{fig:q2gantt}

\subsection{Channel-model sensitivity}
Table~\ref{tab:channel} compares the constant-penalty baseline with the
ITU-R P.~526 single-edge model and its multi-edge upper bound.

\begin{table}[htbp]\centering\small
\caption{Channel-model sensitivity (identical transport plan and candidate
rule; only the occlusion loss differs).}
\label{tab:channel}
\fitwide{%
\begin{tabular}{lccc}
\toprule
 & constant penalty & P.526 single edge & P.526 upper bound\\
\midrule
Outage incidences & \ENchObsInstances & \ENchAltInstances & \ENchAltUbInstances\\
Demand cells & \ENchObsCells & \ENchAltCells & \ENchAltUbCells\\
Usable hover points & \ENchObsCand & \ENchAltCand & \ENchAltUbCand\\
Phases & \ENchObsPhases & \ENchAltPhases & \ENchAltUbPhases\\
$N=3$ feasible? & \ENchObsNC & \ENchAltNC & \ENchAltUbNC\\
Smallest $N$ with a verified schedule & $\ENchObsNStar$ & $\ENchAltNStar$ & $\ENchAltUbNStar$\\
\bottomrule
\end{tabular}}
\end{table}

\noindent Taking the supplied constant penalty as the reference, moving to
knife-edge diffraction \emph{raises} outage demand by $\ENchReqDelta\%$; taking
knife-edge diffraction as the reference instead, the constant penalty admits
$\ENchCandDeltaAbs\%$ \emph{more} usable hover points, i.e.\ the P.~526 supply is
$\ENchCandDeltaAbs\%$ lower. (The two percentages have different bases ---
$1454$ demand instances and $2471$ hover points respectively --- so they are
stated against their own base rather than as one sentence.) Consequently the
smallest fleet for which
we can construct a verified schedule is not invariant across channel conventions:
it is $\ENchObsNStar$ under the supplied constant penalty and $\ENchAltNStar$
under knife-edge diffraction. Two clarifications are needed here, because the
difference matters for how the result may be used.

\emph{First, the negative half is a search result, not a proof.} At $N=3$ under
P.~526 the constructive verifier reaches cell $\ENchAltNReached$ of
$\ENchAltCells$ and stops; no recursion and no integer program was run on the
P.~526 model. We therefore do not claim that three relays are \emph{infeasible}
under P.~526 --- only that we did not find a schedule, at any of the three
candidate budgets tested ($200/400/600$ points), whereas the verifier \emph{did}
find one at $N=4$ ($\ENchAltNFourSorties$ sorties, no uncovered incidence, SOC
minimum $\ENchAltNFourSoc\%$). The constructive verifier is not a monotone
function of $N$ in general --- on one realisation of the controlled terrain
family (terrain translated by $40$ cells) it finds a schedule at $N=3$ and
fails at $N=4$ --- which is a further reason not to read a verifier failure as
evidence of infeasibility. On the reference instance it happens to succeed at
both $N=3$ and $N=4$. The full $4$-relay trajectory is shipped, sortie by sortie, with
positions, departure and link-completion instants, hover duration, energy and
terminal state of charge, so that the positive half can be inspected rather than
taken on trust; it is identical at all three candidate budgets tested.

\emph{Second, the positive half is what carries the operational message.}
Whatever the true minimum under P.~526, the verified schedule needs
$\ENchAltNStar$ relays where the constant-penalty convention needs
$\ENchObsNStar$; the base of the augmentation changes, so a fleet-size claim is a
claim about a \emph{particular} occlusion convention and must be reported together
with it. Figure~\ref{fig:channel} shows the two effects separately and then
compounded: diffraction raises demand and removes usable hover points, and the
smallest constructible fleet is not the same under the two conventions.

\Dfigwide[0.99]{fig_en_channel}{Channel-model sensitivity. Panel (a): outage
demand rises by $\ENchReqDelta\%$ once the constant occlusion penalty is replaced by
knife-edge diffraction. Panel (b): the supply of backhaul-feasible hover points
falls by almost half, because the relay-to-gateway link is diffracted as well.
Panel (c): the two effects compound, and the smallest fleet with a constructed
schedule rises from $3$ to $4$; the dashed line marks the two-relay
inventory.}{fig:channel}

\subsection{Handover premium}
\label{sec:premium}
Table~\ref{tab:baseline} separates the coverage baseline from the executable
model. Three kinds of statement appear in it, and they must not be conflated:
$N_{\mathrm{naive}}$ is the \emph{exact} optimum of the coverage-only relaxation
(Section~\ref{sec:verifier}), supplied by an integer program; $LB_{\mathrm{handover}}$
is a \emph{proven lower bound} for the handover-aware model, obtained as the
minimum number of points covering every union of adjacent phases; and
$N^{*}=3$ is a \emph{constructed} schedule --- an upper bound. Hence
$\max(N_{\mathrm{naive}},LB_{\mathrm{handover}})=\ENrbLB\le N^{*}\le\ENrbGreedy$, and the
single value $\ENrbGreedy$ should be read as ``the smallest fleet for which we
have a verified schedule'', never as the minimum.

\begin{table}[htbp]\centering\small
\caption{Relay-side baseline and bounds on the same instance. The status column
records whether each entry is an exact optimum, a proven bound, or a constructed
schedule.}
\label{tab:baseline}
\fitwide{%
\begin{tabular}{lll}
\toprule
Quantity & Value & Status\\
\midrule
$N_{\mathrm{naive}}$ (instant relocation, per-cell set cover) & \ENrbNaive & exact (IP)\\
$LB_{\mathrm{handover}}$ (adjacent-phase union cover) & \ENrbLB & proven bound (Prop.~\ref{prop:lb})\\
$\ENboundEnergy$ (energy bound, under (H)) & \ENboundEnergy & proven bound\\
Smallest $N$ with a constructed schedule & \ENrbGreedy & constructed\\
Handover premium over the relaxation & \ENrbPremium & ---\\
\bottomrule
\end{tabular}}
\end{table}

\noindent\textbf{Why $LB_{\mathrm{handover}}$ is a lower bound.} The quantity is
computed as the largest, over consecutive phase pairs $(\Phi_i,\Phi_{i+1})$, of
the minimum number of hover points needed to cover \emph{both} phases' demand.
Since a point count is not obviously a fleet count --- one relay may change point
between the two phases --- we record the argument.

\begin{proposition}[Static lower bound from adjacent phases]
\label{prop:lb}
Split the mission into phases at the instants where the set of out-of-range
sorties changes, and let $c_i$ be the minimum number of hover points whose union
covers all demand occurring in phase $\Phi_i$. Then any feasible $N$-relay
schedule satisfies $N\ge \max_i c_i$.
\end{proposition}
\begin{proof}
Fix a feasible schedule and a phase $\Phi_i$. Let $P_i$ be the set of hover
points \emph{actually occupied} during $\Phi_i$ by the $N$ relays. Condition (V)
requires that the union of the points occupied at each instant covers that
instant's demand, and each occupied point is one of the at most $N$ relays. Hence
$P_i$ is a set of at most $N$ points covering all demand in $\Phi_i$, so
$|P_i|\ge c_i$ by definition of $c_i$. Since $|P_i|\le N$, we get $N\ge c_i$, and
$i$ was arbitrary. (Note that this bound does not assume the relays are static
across the phase --- only that at most $N$ points are occupied at any instant, so
a relay that changes point inside a phase only enlarges $P_i$.)
\end{proof}

\noindent Two remarks on the implementation. First, the computed quantity is
$\max_i c_i$ over consecutive phase \emph{pairs}, which is the same bound applied
to the coarser partition obtained by merging adjacent phases; merging can only
lower $\max_i c_i$, so what we report is a valid though looser lower bound,
and on this instance it evaluates to $\ENrbLB$. Second, the bound is genuinely
weaker than the truth: it ignores handover blackout entirely, so it cannot see the
$N=3$ requirement that Theorem~\ref{thm:handover} and the recursion establish. It
is therefore useful as a cheap sanity floor, not as a substitute for the
feasibility argument. Figure~\ref{fig:baseline} displays the three quantities
together and makes the handover premium visible: the coverage abstraction answers
a question nobody is asking, and the executable answer costs
$\ENrbPremium$ more relays.


\Dfigwide[0.94]{fig_en_baseline}{Relay-side baselines. Panel (a): the
coverage abstraction that ignores handovers needs $N_{\mathrm{naive}}=1$ relay,
the rigorous lower bound (minimum covering points of every union of adjacent
phases) is $2$, and the constructive handover-aware verifier finds a schedule
with $3$. The true optimum is bracketed, and the gap between the abstraction and
the executable model --- the \emph{handover premium} --- is $2$ vehicles. Panel
(b): the per-phase-pair lower bound is $2$ throughout, i.e.\ no single phase pair
can be served by one hovering relay.}{fig:baseline}

\Dfigwide[0.98]{fig_en_profile}{Why relaying cannot keep up, in two scales. Panel (a): the six demand phases are $430$--$3190$\,s long, and the $430$\,s short phase is shorter than the median blackout window ($\ENqBlackoutMedian$\,s, dashed) --- a single repositioning costs more than that whole phase. Panel (b): the gap between consecutive phases is $10$\,s throughout, against a cheapest repositioning cost of $\ENwMinMTwo$\,s under M2 and $\ENwMinMOne$\,s under M1. No point change therefore fits into a phase boundary, so any change must occur inside a phase and black out that phase's demand.}{fig:profile}

\paragraph{Why exactly two relays fail, and on what search space.}
The certificate of Corollary~\ref{cor:deficiency} is satisfied at
$\ENqStarBase\%$ of the cells under M1, so it does not close the question, and
Proposition~\ref{prop:energy} only asks for $\ENboundEnergy$. The verdict we
report comes from the recursion of Theorem~\ref{thm:two}, restricted to the
candidate set $\mathcal{C}_0$ of points that can individually cover at least one
demand cell ($\ENcandCZero$ of the $\ENqCand$ hover points) together with the
per-cell top-$12$ refinement: on that set the recursion reports no admissible
state at $t=\ENgrFailNTwoT$\,s under M1 and $t=\ENgrFailNTwoTairPP$\,s under M2, so
no two-relay schedule exists whose hover points lie in that set
(Proposition~\ref{prop:distinguishing}). The pruning scope of that verdict is stated
once in Section~\ref{sec:reproducible}; an unpruned run at the same
breadth was attempted and did not complete within our compute
budget, so the evidence class here is ``no feasible schedule found by the
recursion on a declared finite set'', not an unconditional in-\emph{state-space}
decision. What is unconditional is the inner conclusion: \emph{for points drawn
from $\mathcal{C}_0\cup(\text{top-}12)$}, infeasibility is a statement about a
finite, explicitly listed set and does not depend on the order of expansion.

\noindent\textbf{How the search space itself is chosen.} Because the verdict is
relative to $\mathcal{C}_0$, we state the rule and the scope explicitly.
$\mathcal{C}$ is the dense grid ($0.0035^\circ$ in longitude and latitude,
three hover heights above ground) pre-filtered so that both the access and the
backhaul link can close at some instant, giving $\ENqCand$ points.
$\mathcal{C}_0\subseteq\mathcal{C}$ keeps those points that individually cover at
least one demand cell ($\ENcandCZero$ points), and for each cell we additionally
keep the top-$12$ points by remaining coverage. The two nested statements are
then:
\begin{itemize}[leftmargin=1.4em,itemsep=1pt]
\item \textbf{Unconditional in the set}: if the recursion over
$\mathcal{C}_0\cup(\text{top-}12)$ with pruning disabled finds no admissible
state, no two-relay schedule exists whose hover points lie in that set
(Proposition~\ref{prop:distinguishing}). This is a finite, checkable claim.
\item \textbf{Heuristic in the set}: whether \emph{the instance} has no two-relay
schedule depends on whether the optimal hover points lie in $\mathcal{C}_0\cup
(\text{top-}12)$. We treat that as a modelling assumption, not a theorem. It is
supported by two observations and nothing stronger. First, the relay-side scan of
$\mathcal{C}_0$ itself (the relay-side candidate-set scan shipped in \texttt{results/}) returns no feasible state at
per-cell top-$12$ or top-$24$; the top-$48$ expansion was \emph{not} run to
completion, so it is absent from the evidence rather than negative. Second, the
reported verdict does not change with the candidate budget in the channel
comparison ($\ENchBudgets$ candidates all give the same smallest fleet), which is
weaker evidence, as that budget is the channel experiment's and not the recursion's;
together they are evidence
against --- but not a proof of the absence of --- a better point outside the set.
\end{itemize}
The distinction matters because a referee could otherwise object that the
recursion is being used as an oracle for a set that was itself chosen
heuristically; the items above are the reply.

\noindent We also report where the obstruction lives: in the diagnostics the
rejected transitions are dominated by ``the other relay has not been in place
long enough'', i.e.\ by the alibi requirement, not by the energy budget. Under M1
the handover inequality is satisfiable at the median blackout window for
$\ENqStarBase\%$ of the demand cells under M1 (M2: $\ENqStarAir\%$), so a
two-relay schedule is forced to freeze its
configuration across stretches in which no single configuration covers the
demand --- exactly what the empty state set detects.

\noindent Figure~\ref{fig:theory} shows the three ingredients together: the span
profile $\mathrm{span}(t)$ with its median, the release rate as a function of the
blackout threshold, and the resulting bracket on the fleet size.

\Dfigwide[0.98]{fig_en_theory}{Theory panel. (a) The span profile
$\mathrm{span}(t)$ (top $2\%$ of values shaded by the blue curve) together with
the median blackout windows of both relocation modes; the gaps in which no single
point can hold alone are exactly the stretches in which a repositioning relay
cannot be covered. (b) The fraction of demand cells that admit a handover as a
function of the assumed blackout threshold: at the median M1 window only
$\ENqStarBase\%$ of cells qualify, against $\ENqStarAir\%$ for air transit.
(c) The proven energy bound of Proposition~\ref{prop:energy} set against the
constructive upper bound; the recursion's first empty state is \emph{not} a lower
bound and is not plotted as one --- the interval
$\ENboundEnergy\le N^{*}\le\ENrbGreedy$ is what the instance supports.}{fig:theory}

\subsection{Method cross-validation: a phase-level integer program}
\label{sec:ilp}
Because the negative half of the fleet-size answer rests on one algorithm, we
solve the same decision question a second time with a structurally different
model: a phase-level integer program over the $\ENilpPhases$ demand phases, with
a binary variable per (relay, phase, candidate point), coverage imposed
incidence by incidence, and the blackout condition imposed as an incompatibility
cut whenever $W(a,b)$ exceeds the interval between two consecutive phases. The
model has $\ENilpNTwoBool$ binaries and $\ENilpNTwoCuts$ cuts at $N=2$ and is
solved to proven optimality in $\ENilpNTwoT$\,s. Table~\ref{tab:ilp} collects the
three fleet sizes and places the relaxation's verdict next to the recursion's, so
that the two structurally different procedures can be compared row by row.

\begin{table}[htbp]\centering\small
\caption{Phase-level integer-program relaxation versus the exact cell-level
recursion. The third column records whether CP-SAT \emph{closed} the model
(status \textsc{optimal} or \textsc{infeasible}); note that the shipped result
file's \texttt{proven} field uses the narrower meaning ``proved infeasible'',
so it is \texttt{false} for the two satisfiable rows below. A feasible relaxation
is \emph{not} a schedule (Proposition~\ref{prop:phase}).}
\label{tab:ilp}
\fitwide{%
\begin{tabular}{llll}
\toprule
$N$ & Relaxation status & CP-SAT closed? & Implication\\
\midrule
$1$ & \ENilpNOneStatus & yes ($\ENilpNOneT$\,s) & $N=1$ infeasible, independently confirmed\\
$2$ & \ENilpNTwoStatus & yes & relaxation satisfiable; \emph{not} a schedule\\
$3$ & \ENilpNThreeStatus & yes & relaxation satisfiable (as expected)\\
\bottomrule
\end{tabular}}
\end{table}

\noindent The two models agree on $N=1$ and disagree on $N=2$, and the
disagreement is the interesting part. The relaxation's $N=2$ witness is an
explicit assignment that gives each relay one hover point per phase; what it omits
is that a relay changing its point \emph{inside} a phase leaves a blackout window
that the other relay must cover from a single fixed point. The relevant comparison
is therefore with the \emph{gap between} consecutive phases, which is exactly the
quantity the relaxation's condition (ii) constrains: here every one of the
$\ENilpPhases-1$ gaps is $10$\,s wide, whereas the shortest possible repositioning
costs $w_{\min}=\ENwMinMTwo$\,s under M2 and $\ENwMinMOne$\,s under M1. No
point change can therefore be squeezed into a phase boundary at all, so any change
must happen inside a phase and black out that phase's demand. The speculatively
relevant alternative --- phase \emph{lengths}, which here range from $430$\,s to
$3190$\,s --- is not the binding quantity, and is quoted only to show that the
gap, not the length, is what a relocation must fit into. The span profile
($\ENspanMedian$\,s median, $\ENspanMax$\,s maximum) is likewise comparable to the
blackout window ($\ENqBlackoutMedian$\,s median) rather than to a phase. The
bottleneck is therefore
not the phase-granular resource count --- which the relaxation shows is
comfortably satisfied at $N=2$ --- but the intra-phase temporal coupling that only
the cell-level recursion can see. This is the sharpest form of the answer to
``is two relays enough'': \emph{satisfiable at phase granularity, no schedule
found at instant granularity}, and the gap is quantified by the release rates of
Figure~\ref{fig:theory}.

\noindent The relaxation is solved on a reduced candidate set (the
$\ENilpPerPhase$ points per phase that cover the most incidences), so its
\emph{positive} answers carry that restriction, exactly as the recursion's do;
its negative answer for $N=1$ is strong for the same reason and agrees with both
other methods ($N_{\mathrm{naive}}=1$ coverage abstraction is necessary but not
sufficient; the constructive verifier fails at $t=\ENgrFailNOneT$\,s).

\Dfigwide[0.98]{fig_en_methods}{Three independent decision procedures on the same instance. Each cell is the verdict of one procedure at one fleet size: the state recursion of Theorem~\ref{thm:two} (M1), the phase-level integer program of Section~\ref{sec:ilp}, and the constructive verifier of Section~\ref{sec:verifier}. The three agree on $N=1$ (no schedule) and, between the two procedures that were run, on the existence of a schedule at $N=3$ and $N=4$; they disagree only at $N=2$, where the relaxation is satisfiable while the instant-level recursion finds no state. That disagreement is the relaxation gap of Proposition~\ref{prop:phase}, and it is why the relaxation is reported alongside the recursion rather than instead of it. The recursion was not run for $N\geq3$; those cells are left blank rather than filled with an assumption.}{fig:methods}

\subsection{Relocation mode: depot return versus air transit}
A natural objection is that infeasibility is an artefact of the depot-return
assumption: if a relay could fly directly from one hover point to another, the
blackout window would shrink and two relays might suffice. We therefore
implemented and compared both modes under the same candidate set and the same
energy budget.

\Dfigwide[0.90]{fig_en_modes}{Relocation modes. Panel (a): air transit cuts the
median blackout window from $\ENqBlackoutMedian$\,s to $\ENqBlackoutAirMedian$\,s. Panel (b): the necessary
condition of Theorem~\ref{thm:handover} is satisfied on $\ENqStarBase\%$ of the demand
cells under M1 and on $\ENqStarAir\%$ under M2 --- a real improvement, but the condition
still fails somewhere, and no $N=2$ schedule was found under \emph{either} mode
(pruned recursion over the declared set).}{fig:modes}

\noindent The improvement is genuine in magnitude but does not change the
conclusion: Figure~\ref{fig:modes} shows the two effects, and with two relays the
obstruction survives both relocation models, so the result is not a consequence of
the depot-return assumption. Table~\ref{tab:modes} gives the verdict fleet size by
fleet size.

\begin{table}[htbp]\centering\small
\caption{Dual-mode feasibility verdict. The source column distinguishes the two
procedures: \emph{recursion} is the state recursion of
Theorem~\ref{thm:two} on the declared candidate set with a per-cell top-$12$
refinement (run with its reported pruning policy, not unpruned), and
\emph{verifier} is the constructive height-greedy procedure. A
negative verdict from either is a search failure over its scope, not a proof of
infeasibility, and the two procedures are not interchangeable: the verifier is
not monotone in $N$ in general (on a translated-terrain realisation of the
controlled family it finds a schedule at $N=3$ and fails at $N=4$), so its
failures carry less weight than the recursion's.}
\label{tab:modes}
\setlength{\tabcolsep}{3pt}
\fitwide{%
\begin{tabular}{@{}p{0.4cm}p{6.0cm}p{6.0cm}p{1.6cm}@{}}
\toprule
$N$ & Mode M1 (depot return) & Mode M2 (air transit) & Source\\
\midrule
$1$ & no state, $t=\ENgrFailNOneT$\,s & no state, $t=\ENgrFailNOneT$\,s & recursion\\
$2$ & no state, $t=\ENgrFailNTwoT$\,s (cell \ENgrFailNTwoTCell) &
no state, $t=\ENgrFailNTwoTairPP$\,s (cell \ENgrFailNTwoTairPPCell) & recursion\\
$3$ & \textbf{feasible}: \ENgrNThreeMOneSorties{} sorties,
\ENgrNThreeMOneMiss{}/\ENqInstances{} uncovered, SOC min
$\ENgrNThreeMOneSoc\%$ & no schedule found: reached cell
\ENgrNThreeMTwoReached{} of \ENqCells{} at $t=\ENgrNThreeMTwoFailT$\,s,
after \ENgrNThreeMTwoActions{} relocations & verifier\\
\bottomrule
\end{tabular}}
\end{table}

\noindent Three readings of Table~\ref{tab:modes} matter for the reviewer's
objection. First, the objection is answered: $N=2$ fails under \emph{both}
relocation models, and under M2 the recursion survives to
$t=\ENgrFailNTwoTairPP$\,s --- $3.3$ times further in \emph{elapsed time}
(cell index $684$ versus $171$, a ratio of $4.0$) than the $t=\ENgrFailNTwoT$\,s
reached under M1 --- before the state set empties, so air
transit does buy progress, just not enough. Second, at $N=3$ the two modes
diverge in our search: M1 is certified by an explicit zero-miss schedule, while
M2 is not certified. Third, that divergence must be read carefully. The M2
verifier made \ENgrNThreeMTwoActions{} relocations before stalling at
$t=\ENgrNThreeMTwoFailT$\,s, against eight in the M1 run that succeeded: the
air-transit heuristic spends its energy budget on more, shorter hops, and the
schedule is then cut short by the airborne-time cap rather than by geometry. We
therefore report M2 at $N=3$ as ``no schedule found by the verifier'', not as
infeasible, and we do \emph{not} claim an air-transit fleet saving: the honest
statement is that air transit reduces the blackout window by a factor
$\ENqBlackoutMedian/\ENqBlackoutAirMedian$ without changing the $N=2$ verdict, so
the conclusion that the partition needs three relays is robust to
Assumption~\ref{asm:handover}'s depot-return clause.

\subsection{Is decoupling the cause?}
Table~\ref{tab:feedback} reports the co-design loop, in which only the departure
instants of the transport sorties are varied; batching, routes, aircraft types
and assignments are frozen. The deficiency of
Corollary~\ref{cor:deficiency} was unchanged by that search: the vector returned
at the end of the search is the baseline vector itself (zero sorties shifted, and
the trace contains no strictly improving move). Note that the acceptance test is
threshold-accepting, so equal-deficiency shifts may be walked through during the
search; what is reported is that none of them survives as the returned incumbent --- so on
this frozen plan, searching over departure instants alone did \emph{not} reduce
the obstruction. We state this as what it is: a null result on one coordinate of
the design space. Because routing, batching, aircraft assignment and fleet
composition were all held fixed, it is evidence against a purely
scheduling-driven explanation of the residual deficiency, not a proof that
decoupling is harmless, and not a statement about the joint design space.

\begin{table}[htbp]\centering\small
\caption{Co-design feedback loop over departure instants only. The deficiency is
counted in demand cells; demand cells are not uniformly spaced in time, so no
seconds conversion is reported.}
\label{tab:feedback}
\fitwide{%
\begin{tabular}{lcccc}
\toprule
Mode & deficiency before & after & reduction & zero?\\
\midrule
M1 (depot return) & \ENcfBaseDefA & \ENcfBaseDefB & $\ENcfBaseRed\%$ & \ENcfBaseZero\\
M2 (air transit) & \ENcfAirDefA & \ENcfAirDefB & $\ENcfAirRed\%$ & \ENcfAirZero\\
\bottomrule
\end{tabular}}
\end{table}

\Dfigwide[0.94]{fig_en_feedback}{Co-design feedback loop. Panel (a): the
handover deficiency is identical before and after re-optimising the departure
instants, and the returned vector is the baseline vector (no shift survives as the
incumbent), so on this frozen plan the
departure instants are not where the residual deficiency can be removed. Panel
(b): the span profile --- the longest interval over which one single hover point
can carry the whole demand --- repeatedly falls below the required blackout
window (dashed lines), which is the mechanism behind the obstruction.}{fig:feedback}

\subsection{Baselines and ablations}
\label{sec:ablation}
Table~\ref{tab:ablation} collects the transport-side method comparison and the
partitioning cost. Two caveats apply to that table. The method comparison runs under
a fixed iteration budget that differs from the budget used for the frozen plan of
Section~\ref{sec:exp}, so its sortie counts are comparable across \emph{methods}
but not against the frozen plan; and the rows drawn from the two weight sets
(sortie-minimising and on-time-first) are not interchangeable, which is why the
CP-SAT status row is drawn from the on-time-first set while the two sortie counts
are drawn from the sortie-minimising one. On the transport layer, a pure construction heuristic needs
$\ENbaseGreedySorties$ sorties and $\ENbaseGreedyEnergy$\,kWh on the
sortie-minimising weight set, against $\ENbaseAlnsSorties$ sorties and
$\ENbaseAlnsEnergy$\,kWh for ALNS; the CP-SAT polishing stage returns an
\ENbaseCpsatStatus{} schedule on its neighbourhood. Partitioning the service
area into $K$ autonomous groups raises the weighted resource demand from
$\ENpopKtwo$ ($K=2$) to $\ENpopKthree$ ($K=3$), the \emph{price of
partitioning}: splitting duplicates per-group peaks and therefore inflates the
fleet.

\paragraph{A textbook baseline, and what it gets wrong.}
The four transport methods above all belong to our own algorithmic family, which
invites the objection that the comparison is self-referential. We therefore added
the most widely cited constructive heuristic for this problem class, the
Clarke--Wright savings algorithm \cite{clarke1964}, generalised from endpoint
merging to multi-stop tours as
$\mathrm{save}(A,B)=t(A)+t(B)-t(A\cup B)$ on sortie durations and applied until
no capacity-, volume- and energy-feasible merge remains. It is evaluated with the
\emph{same} scheduler and the same metric function as ALNS, so the comparison is
apples-to-apples (Table~\ref{tab:cw}).

\begin{table}[htbp]\centering\small
\caption{Clarke--Wright savings versus ALNS under identical scheduling and
evaluation code, on the sortie-minimising weight set. The savings heuristic
dominates the routing-layer objectives and fails the scheduling ones.}
\label{tab:cw}
\fitwide{%
\begin{tabular}{lrrrrr}
\toprule
Method & Sorties & Energy (kWh) & Tardiness & Hard viol. (s) & On-time\\
\midrule
Clarke--Wright & $\ENcwSorties$ & $\ENcwEnergy$ & $\ENcwTardy$ & $\ENcwHard$ & $\ENcwOntime\%$\\
ALNS & $\ENcwAlnsSorties$ & $\ENcwAlnsEnergy$ & $\ENcwAlnsTardy$ & $\ENcwAlnsHard$ & $\ENcwAlnsOntime\%$\\
\bottomrule
\end{tabular}}
\end{table}

\noindent The savings heuristic starts from $\ENcwStart$ single-box sorties
($\ENcwStart$ vehicles) and merges them down to $\ENcwSorties$ in
$\ENcwMerges$ merges, reaching $\ENcwEnergy$\,kWh against ALNS's
$\ENcwAlnsEnergy$\,kWh: it wins decisively on the routing layer. It loses just as
decisively on the scheduling layer, with $\ENcwHard$\,s of hard-deadline
violation against $\ENcwAlnsHard$\,s and an on-time rate of $\ENcwOntime\%$
against $\ENcwAlnsOntime\%$. The mechanism is structural rather than incidental:
a route-level merge creates large sorties that must be dispatched early, and they
then block each other on the shared drone-plus-battery timeline. A solution with
$\ENcwHard$\,s of hard-deadline violation is not refused by our objective ---
hard deadlines enter as a penalty, not as a hard constraint --- but it is charged
every second, so the objective that admits it is $\ENcwAlnsHard$\,s worse on that
term alone. This is the concrete
argument for keeping scheduling inside the search loop instead of repairing it
afterwards --- the same conclusion that the relay side reaches from the opposite
direction (a coverage abstraction that ignores handovers is not executable).

\begin{table}[htbp]\centering\small
\caption{Baselines and ablations (all values from the shipped result files).}
\label{tab:ablation}
\fitwide{%
\begin{tabular}{lll}
\toprule
Layer & Configuration & Outcome\\
\midrule
Transport & pure construction & $\ENbaseGreedySorties$ sorties, $\ENbaseGreedyEnergy$\,kWh\\
Transport & Clarke--Wright savings \cite{clarke1964} & $\ENcwSorties$ sorties, $\ENcwEnergy$\,kWh, $\ENcwHard$\,s hard violation\\
Transport & ALNS w/o CP-SAT (ablation) & $\ENbaseAlnsSorties$ sorties, $\ENbaseAlnsEnergy$\,kWh\\
Transport & ALNS $+$ CP-SAT & \ENbaseCpsatStatus{} on the fixed neighbourhood\\
Relay & coverage only (instant relocation) & $N_{\mathrm{naive}}=\ENrbNaive$\\
Relay & handover-aware, constructive & $N^{*}=\ENrbGreedy$\\
Modes & M1 vs M2 blackout median & $\ENqBlackoutMedian$\,s vs $\ENqBlackoutAirMedian$\,s\\
Modes & M2 $(\star)$ rate & $\ENqStarAir\%$ (vs $\ENqStarBase\%$ for M1)\\
Relay & channel model & $N^{*}=\ENchObsNStar$ (constant) vs $\ENchAltNStar$ (P.526)\\
Coupling & frozen schedule vs co-designed & deficiency unchanged\\
Partitioning & $K=1/2/3$ weighted demand & $1.00$ / $\ENpopKtwo$ / $\ENpopKthree$\\
\bottomrule
\end{tabular}}
\end{table}

\subsection{What actually determines the fleet size}
Figure~\ref{fig:terrain} and Table~\ref{tab:terrain} report the controlled
terrain family. After correcting a node-elevation inconsistency (attachment
elevations belong to the original terrain and must be re-derived once the
elevation field changes; otherwise the aircraft is routed below ground and the
resulting outages are spurious), the fleet size turns out to be non-monotone in
terrain severity and, among the predictors we computed, no more than weakly
associated with the temporal fragmentation of the outage demand. Terrain
realisations with almost identical occlusion ($44.4\%$, $44.6\%$, $45.2\%$) yield
$N^{*}=3,3,>4$ respectively, while a realisation with only $13.1\%$ occlusion
yields $N^{*}>4$: three scenarios with the same fragmentation count
($0$ phases shorter than $1400$\,s) yield $N^{*}=\,>4,\,3,\,2$. That spread is the
honest summary of this experiment --- the pattern is non-monotone, and with nine
scenarios we cannot rank the candidate predictors.

\begin{table}[htbp]\centering\small
\caption{Controlled terrain family. $N^{*}>4$ denotes ``no feasible schedule
found for $N\le 4$'', i.e.\ a censored observation, not a measured value;
$\ENtbNGtFour$ of the $\ENtbNOk$ realisations are censored. The correlations are
Pearson coefficients with the censored realisations entered at an arbitrary code
$5$; the last column gives $95\%$ confidence intervals by the Fisher $z$-transform, all of which
span zero.}
\label{tab:terrain}
\fitwide{%
\begin{tabular}{lccc}
\toprule
Metric & Range & $\rho$ with $N^{*}$ & $95\%$ CI \\
\midrule
Occlusion ratio & $\ENtbOccMin$--$\ENtbOccMax$\% & $-0.23$ & $[-0.77,+0.52]$\\
Phases shorter than $1400$\,s & -- & $+0.47$ & $[-0.28,+0.86]$\\
Outage incidences & -- & $-0.12$ & $[-0.73,+0.59]$\\
\bottomrule
\end{tabular}}
\end{table}

\Dfigwide[0.98]{fig_en_terrain}{Controlled terrain family (only the DEM
elevation field is transformed; nodes and all parameters are frozen, and the
node elevations are re-aligned so that no settlement ends up below ground).
Panel (a): the smallest fleet with a constructed schedule is \emph{not} monotone
in the occlusion ratio --- it rises to $>3$ at the low end, drops to $3$ and then
to $2$ in the middle, and rises again past $45\%$. Panel (b): the
outage demand itself grows monotonically with the occlusion ratio, so panel (a)
is \emph{not} explained by demand volume. Nine scenarios are too few to establish
which predictor ranks first; the largest coefficient we computed is for temporal
fragmentation (Table~\ref{tab:terrain}, $\rho=+0.47$, CI spans
zero).}{fig:terrain}

Figure~\ref{fig:terrain} plots both panels: the non-monotone step function of
panel (a), and panel (b), which rules out the obvious explanation that demand
volume alone drives the fleet size.

\subsection{Price of partitioning}
\label{sec:partition}
The integrated ($K=1$) resource peak is a lower bound on the sum of the
per-group peaks; the gap quantifies the cost of autonomous operation. Under the
inventory-feasible two-relay plan the per-group relay demand is
\ENqfourAKtwoRelayNeed{} (integrated) versus \ENqfourBKtwoRelayNeed{} under the
augmented, continuously-connected plan for $K=2$; with the channel model of
Section~\ref{sec:channel} the same accounting leaves a residual relay gap.

%==============================================================================
\section{Discussion}
\label{sec:discussion}
\paragraph{Relay sizing needs an executable model.}
The coverage abstraction answers ``how many relays must be in the air at the
worst instant''; the operator needs ``how many relays must exist'', and the two
differ by the handover premium ($\ENrbPremium$ vehicles here). Theorem~\ref{thm:handover}
gives a checkable necessary condition that can be evaluated from the demand
profile and the blackout function alone, without solving the scheduling problem;
it is not a sufficient certificate, and we quantify how far it is from closing the
question ($\ENqStarBase\%$ of the demand cells under M1). Theorem~\ref{thm:two} states that
the recursion \emph{would be} exhaustive over a stated candidate set with pruning
disabled; the runs we report use a per-cell top-$12$ refinement, so the two-relay
verdict is a statement about that declared search space --- finite and explicitly
listed --- rather than about the instance. Proposition~\ref{prop:energy} supplies
an independent, sampling-free lower bound under hypothesis (H), and the
constructive verifier a matching upper bound, so the fleet-size interval
$\ENboundEnergy\le N^{*}\le\ENrbGreedy$ is what the instance supports.

\paragraph{The channel convention is part of the answer.}
Replacing the constant occlusion penalty with ITU-R P.~526 knife-edge
diffraction changes both the demand ($\ENchReqDelta\%$) and the supply
($\ENchCandDeltaAbs\%$ of usable hover points), and raises the smallest fleet for
which we can construct a verified schedule from $\ENchObsNStar$ to
$\ENchAltNStar$. We do \emph{not} claim that three relays are infeasible under
P.~526: only a constructive verifier was run there, and it found a schedule at
$\ENchAltNStar$ but not at $3$. Fleet-size statements are therefore statements
about a channel convention \emph{and} a search procedure, and must be reported
with both.

\paragraph{Departure instants are not where the obstruction lives.}
On the frozen baseline plan, re-optimising departure instants left the obstruction
unchanged, so on that plan better delivery scheduling does not close the
communication gap; the gap must be managed operationally (pre-programmed
autonomous legs, a risk register) or removed by adding relays. Because routing,
batching and fleet composition were held fixed, this is a statement about one
coordinate of the design space and not about co-design in general.

\paragraph{Occlusion volume is a poor guide, but nine scenarios prove little.}
In the terrain study occlusion volume does not order the fleet size, and the
non-monotone pattern is real in the sense that realisations with the same
fragmentation count take three different values. With $\ENtbNOk$ scenarios, of
which $\ENtbNGtFour$ are censored at $N^{*}>4$, none of the candidate predictors
shows an association whose confidence interval excludes zero, so we present this
as a caution against extrapolating fleet size from a single ``blocked fraction''
rather than as a ranking of predictors.

%==============================================================================
\section{Reproducibility and limits of the reported evidence}
\label{sec:reproducible}
\paragraph{What is generated, and what is quoted.}
Every table entry and every number in the prose is written by
\texttt{code/make\_numbers\_en.py} into a macro file
(\texttt{numbers\_en.tex}, $\ENnMacros$ macros), so that no reported value is
typed by hand; each macro is traced to its source file in the audit file the
script emits. The exclusions are, in full: the correlation coefficients and
confidence intervals of Table~\ref{tab:terrain} and the terrain occlusion
percentages, computed in the terrain report; the mechanism parameters
($\Delta=10$\,s, the $1400$\,s short-phase cut-off, the per-cell top-$12$ depth,
the $200/400/600$ candidate budgets) together with directly-read summary
statistics such as the phase-length range $430$--$3190$\,s; the $95\%$ level and
censoring code $5$ of Table~\ref{tab:terrain}; and a few once-reported values
(the unpruned $1501$\,s/$555$\,MB budget, the $K=1$ baseline of
Table~\ref{tab:ablation}, the $3.3$ time ratio, the relocation counts $8$ and
$\ENgrNThreeMTwoActions$). Each of these is quoted verbatim from the result file
named beside it.

Every other figure and table value is regenerated from the shipped result files by
the pipeline below, and each macro is traced to its source file in
\texttt{numbers\_en\_sources.txt}. Those file names are Chinese; we keep non-ASCII
characters out of the manuscript source for portability.

\paragraph{Three evidence classes.}
We use three labels throughout, and they are not interchangeable:
\begin{itemize}[leftmargin=1.4em,itemsep=1pt]
\item \textbf{Proven}: a mathematical statement with a proof in the text. These
are Theorem~\ref{thm:handover}, Corollary~\ref{cor:deficiency},
Theorem~\ref{thm:exact}, Theorem~\ref{thm:one},
Proposition~\ref{prop:phase}, Proposition~\ref{prop:energy} (conditional on
hypothesis (H)), Proposition~\ref{prop:lb},
Proposition~\ref{prop:W}, Proposition~\ref{prop:monotone},
Proposition~\ref{prop:demand}, Proposition~\ref{prop:lipschitz},
Proposition~\ref{prop:distinguishing}, Proposition~\ref{prop:structural}, and the
lower bound $LB_{\mathrm{handover}}=\ENrbLB$, which is the exact optimum of a set-cover
integer program over unions of adjacent phase pairs.
\item \textbf{Constructed}: an explicit, independently re-verified schedule.
The $N=\ENrbGreedy$ relay plan and the transport plan are in this class; we
re-check coverage, timing, energy and state of charge with two independent code
paths before reporting them.
\item \textbf{Search scope}: ``no feasible schedule was found'' by a specified
procedure on a specified finite set. Every negative fleet-size verdict in this
paper is in this class --- including $N=2$. In particular the recursion's verdict
is over the declared candidate set $C_0$ with a per-cell top-$12$ refinement and
the pruning policy reported in Section~\ref{sec:reproducible}; the run reported
here is the pruned one, and the unpruned run is discussed below.
We therefore never write ``proven impossible'' for a negative fleet-size verdict.
\end{itemize}

\paragraph{What we did not manage to establish.}
We record these explicitly so that the scope of the claims is unambiguous.
(i) We could not decide the two-relay question over the full candidate set, and we
could not even re-run it with state pruning disabled at the smaller top-$12$
breadth: an unpruned attempt consumed $1501$\,s of CPU (peak $555$\,MB) without
terminating and was stopped by hand, so the verdict stands as a
\emph{pruned} recursion over $C_0$ with the top-$12$ refinement. That negative
result is itself shipped as a result file, with the exact invocation, so a reader
can reproduce the non-termination rather than take our word for it. (ii) We could
not turn the span inequality into a search-free certificate. Made parameter-free
---
$W(a,b)\le\mathrm{span}_{\min}(k)$ with no tuned threshold --- it holds for at
least one pair in $\ENpfCellsOkBase$ of $\ENpfCells$ cells under M1 and
$\ENpfCellsOkAir$ of $\ENpfCells$ under M2, so it cannot certify infeasibility.
(iii) Under the P.~526 channel model no exact decision procedure was run at all;
only the constructive verifier was, and it is non-monotone in $N$. (iv) The
terrain study rests on $\ENtbNOk$ scenarios with $\ENtbNGtFour$ censored, which is
too few to rank predictors, and the ``phase transition'' reading is descriptive
rather than tested for a bifurcation. (v) The demand profile is evaluated on a
fixed sampling grid, so outage statistics are grid-verified rather than proved in
continuous time; Corollary~\ref{cor:continuous} shows the qualitative conclusion
survives the continuous limit, but the counts do not.

%==============================================================================
\section{Conclusion}
\label{sec:conclusion}
We recast hovering-relay feasibility as a handover-aware scheduling problem,
proved a checkable necessary condition and a span-based obstruction, replaced a
constant occlusion penalty by ITU-R P.~526 knife-edge diffraction, ran a
departure-instant-only feedback loop, and gave a nine-scenario sensitivity study
of how the fleet size varies across a controlled terrain family.
Limitations are stated explicitly and are not incidental. The two-relay verdict,
like every other negative fleet-size verdict here, is a statement about a
declared finite search space --- the candidate set $C_0$ with a per-cell top-$12$
refinement, evaluated with the pruning policy we report --- and never ``proven
impossible''; closing state pruning at that breadth did not terminate within our
budget, so the verdict is relative to that declared space. The energy bound
(Proposition~\ref{prop:energy}) is sampling-free but conditional on hypothesis
(H), and Proposition~\ref{prop:structural} and the span inequality are necessary
conditions, not certificates. Under the P.~526 model only a constructive verifier
was run, so we claim a constructed schedule at $\ENchAltNStar$ relays and not
infeasibility at three. For $N\ge3$ the verifier is a sufficient procedure, so its
failures are search failures. The transport layer is solved heuristically, so no
optimality gap is claimed: the reported transport figures are the best feasible
solutions found. The terrain family is derived from a single elevation model by
controlled transformations, and the outage statistics are verified on a fixed
sampling grid rather than in continuous time. Finally, the reference instance uses
the elevations supplied with the benchmark; the terrain-family study re-derives
them from the transformed elevation field.

% ---- 声明区：Elsevier 要求投稿正文自带这些声明（不能只放在单独文件里）----
% 先前只放在投稿包的 05/06 单独文本里，正文没有；Elsevier 的 Guide for Authors
% 要求 Data availability 等声明出现在稿件中，故补在此处（结论之后、参考文献之前）。
\section*{Declarations}
\addcontentsline{toc}{section}{Declarations}

\noindent\textbf{Data availability.}
The instance (node coordinates, cargo manifest, fleet and radio parameters) is the
benchmark shipped with the problem statement. The 30\,m digital elevation model is
Copernicus GLO-30 (ESA), openly available. All derived result files, the figure and
table generators, and a one-command reproduction pipeline are archived at
\url{https://github.com/ccx111145/data}. Every numerical value in this manuscript
is regenerated from those files by \texttt{code/make\_numbers\_en.py}; the
bibliography was checked entry by entry against its publisher record by DOI.

\noindent\textbf{CRediT authorship contribution statement.}
C.~Cao: Conceptualization, Methodology, Software, Formal analysis,
Writing -- original draft. Q.~Li: Methodology, Software, Validation,
Writing -- review \& editing. Z.~Liu: Conceptualization, Methodology,
Formal analysis, Supervision, Project administration.
Y.~Peng: Investigation, Data curation, Visualization.
Z.~Ouyang: Investigation, Resources, Data curation, Writing -- review \& editing,
Visualization.

\noindent\textbf{Declaration of competing interest.}
The authors declare that they have no known competing financial interests or
personal relationships that could have appeared to influence the work reported in
this paper.

\noindent\textbf{Acknowledgements.}
None.

\noindent\textbf{Funding.}
This research received no external funding.

% ---- 编译顺序：xelatex -> bibtex -> xelatex -> xelatex -------------------
\bibliographystyle{unsrt}
\bibliography{refs}

\end{document}
